\subsubsection{圆维签名与零色散二阶侧信息：结构圆因子与残差账本的可辨识接口}\label{subsubsec:xi-cdim-signature-zero-dispersion}
第 \ref{sec:circle-dimension-phase-gate} 节以 Pontryagin 对偶的连通分量维数定义圆维，并在 \S\ref{subsec:cdim-implementation-audit} 将不可逆读出制度化为外置残差寄存器以维持账本闭合。
另一方面，\S\ref{subsubsec:xi-time-protocol-conclusions}--\S\ref{subsubsec:xi-time-protocol-conclusions-part2} 的时间协议结论链在零色散口径下给出了最小侧信息预算 \(L_m(\varepsilon)=m\gamma+b(\varepsilon)+o(1)\) 的常数项反演（推论 \ref{cor:xi-constant-bias-inversion-b-eps}）。
为在结构层与实现层之间建立可组合的审计接口，本小节引入把连续圆因子与离散残差轴分离记录的圆维签名，并给出有限原子零色散下常数项曲线的闭式、折点结构、可识别性与串接律。
\par

\begin{definition}[Circle-dimension signature / 圆维签名 \texorpdfstring{$\Sigma\Sigma$}{SigmaSigma}]
\label{def:cdim-signature}
沿用定义 \ref{def:circle-dimension}--\ref{def:half-circle-dimension} 的圆维 \(\cdim\) 与半圆维 \(\cdim_{+}\)。
设 \(X\) 为一个\emph{可审计信息对象}，并假定它在结构层至少携带：
\begin{enumerate}
  \item 一个有限生成离散交换群 \(G_X\)（可逆相位通道/可见交换结构）；
  \item （可选）一个有限生成可消去交换幺半群 \(M_X\)，作为 \(G_X\) 的正锥/单侧约束；
        若 \(M_X\) 给定，则其 Grothendieck 群完成记为 \(G(M_X)\)。
\end{enumerate}
在实现层，约定任何不可逆读出均按扩展保存口径外置为残差寄存器（命题 \ref{prop:unit-circle-phase-extension-preserving}，并见 \S\ref{subsec:cdim-implementation-audit}）。
\smallskip

\noindent\textbf{结构圆维剖面（circle profile）.}
定义
\[
\Sigma_{\circ}(X)
:=\bigl(\cdim(G_X),\,\cdim_{+}(M_X)\bigr)
\in (\NN\cup\{0\})\times(\tfrac12\NN\cup\{0\}),
\]
其中当 \(M_X\) 不给定时约定 \(\cdim_{+}(M_X)=0\)。
\smallskip

\noindent\textbf{残差账本签名（ledger / register signature）.}
定义 \(\Sigma_{\mathrm{res}}(X)\) 为一组\emph{外置记录轴/残差寄存器}的有限元组，
用于承载一切不产生新圆因子的“纤维标签/细分/截断/量化”等实现信息。
在单位圆相位门语境下，可取规范三元模板
\[
\Sigma_{\mathrm{res}}(X)
:=\bigl(\mathsf{End}_\varepsilon(X),\,\mathsf{Prec}(X),\,\mathsf{Tag}(X)\bigr),
\]
分别记录端点窗口质量、精度势累计以及不可寻址通道的类型标签（\S\ref{subsec:cdim-implementation-audit}）。
此外，任何有限（或由有限生成 torsion 群生成的）纤维扩张与离散标签插入均并入 \(\mathsf{Tag}(X)\)，
并在结构层不改变圆维（命题 \ref{prop:circle-dimension-laws}）。
\smallskip

\noindent\textbf{圆维签名（double-sigma）.}
定义 \(X\) 的\emph{圆维签名}为
\[
\Sigma\Sigma(X)
:=\bigl(\Sigma_{\circ}(X),\,\Sigma_{\mathrm{res}}(X)\bigr).
\]
因此，\(\Sigma\Sigma\) 的第一分量度量“独立相位圆通道数”，第二分量保存所有
\emph{圆维为零但在实现/审计中不可忽略}的额外信息（账本/寄存器/纤维标签）。
\end{definition}

\begin{proposition}[Implementation-layer coding / 实现层编码不导致结构圆维坍缩]
\label{prop:impl-coding}
固定有限字母表 \(\mathcal{A}\)，令 \(\mathcal{A}^\ast\) 为所有有限字符串的集合。
存在一个可计算单射编码
\[
\mathrm{enc}:\mathcal{A}^\ast \hookrightarrow \NN.
\]
因此：只要 \(X\) 及其证书/账本 \(\Sigma_{\mathrm{res}}(X)\) 具有有限描述
\(\mathrm{desc}(X)\in\mathcal{A}^\ast\)，就可以在\emph{实现层}用
\[
\mathrm{code}(X):=\mathrm{enc}\bigl(\mathrm{desc}(X)\bigr)\in\NN
\]
进行无损序列化存储，并可由解码器恢复该描述。
\par
该结论仅陈述“实现层可单射化/可存储”，不蕴含任何“结构圆维坍缩”：
\begin{enumerate}
  \item \(\cdim(G_X)\) 与 \(\Sigma\Sigma(X)\) 是 \(X\) 的\emph{结构层}不变量（按同构类计），
        不应由 \(\mathrm{code}(X)\) 的算术结构反演得到；
  \item 编码 \(\mathrm{code}\) 一般不是群同态且不保持运算，因此不用于比较结构圆维；
  \item 在本文语义中，这正是“扩展保存”的一个实现层实例：当某读出 \(\pi\) 非单射时，
        通过外置记录轴 \(r\) 使扩展读出 \((\pi,r)\) 单射，从而把不可逆性显式移入账本/寄存器
        （命题 \ref{prop:unit-circle-phase-extension-preserving}）。
\end{enumerate}
\end{proposition}
\leanverified{paper\_impl\_coding}
\begin{proof}
可取任意固定的前缀无歧义编码并将字符串解释为自然数（例如以 \(|\mathcal A|\)-进制表示并追加分隔符），即得到可计算单射 \(\mathrm{enc}\)。
其余断言由定义解释立得。证毕。
\end{proof}

\begin{theorem}[圆维缺口的条件 Kolmogorov 不可消去性：相位通道不足时残差典型不可压缩]\label{thm:xi-cdim-kolm-residual-gap}
设 \(G\cong \ZZ^{r}\oplus T\) 为有限生成交换群，其中 \(T\) 有限。
对每个精度 \(b\ge 1\) 定义窗口
\[
W_b(G):=\{0,1,\dots,2^b-1\}^{r}\times T,
\qquad
|W_b(G)|=2^{br}|T|.
\]
令相位字母表为 \(\mathrm{Phase}(k)_b:=(\ZZ/2^b\ZZ)^k\)。
设存在一族可计算单射
\[
\mathrm{Enc}_b:W_b(G)\hookrightarrow \mathrm{Phase}(k)_b\times R_b,
\]
其中 \(R_b\) 为有限残差寄存器集合。
记 \(\mathrm{Enc}_b(x)=(\pi_b(x),\rho_b(x))\) 为相位与残差分量。
固定一台前缀自由通用机，并以 \(K(\cdot\mid\cdot)\) 表示条件 Kolmogorov 复杂度。
则对任意 \(t\ge 1\)，存在常数 \(C>0\)（与 \(b,t\) 无关，仅依赖所选通用机与 \(\mathrm{Enc}_b\) 的实现范式）使得
\[
\card{\Bigl\{x\in W_b(G):\ K\bigl(\rho_b(x)\mid \pi_b(x),b\bigr)
<
(r-k)b+\log_2|T|-t-C\log b\Bigr\}}
\ \le\ 2^{-t}\,|W_b(G)|.
\]
特别地，若 \(X\sim\mathrm{Unif}(W_b(G))\)，则
\[
\PP\Bigl[
K\bigl(\rho_b(X)\mid \pi_b(X),b\bigr)
\ge
(r-k)b+\log_2|T|-t-C\log b
\Bigr]\ \ge\ 1-2^{-t}.
\]
\end{theorem}
\leanverified{paper\_xi\_cdim\_kolm\_residual\_gap}
\begin{proof}
记 \(N_b:=|W_b(G)|=2^{br}|T|\)，并取
\[
m:=\log_2 N_b-t-C_0
=br+\log_2|T|-t-C_0.
\]
由 Kolmogorov 复杂度的计数性质，满足 \(K(x\mid b)<m\) 的 \(x\in W_b(G)\) 至多 \(<2^m=2^{-t}N_b\) 个。
因此除至多 \(2^{-t}N_b\) 个例外外，都有
\begin{equation}\label{eq:xi-kolm-residual-gap-lb-Kx}
K(x\mid b)\ \ge\ br+\log_2|T|-t-C_0.
\end{equation}
对任意满足 \eqref{eq:xi-kolm-residual-gap-lb-Kx} 的 \(x\)，由于 \(\mathrm{Enc}_b\) 可计算且为单射，有
\[
K(x\mid b)\le K\bigl(\pi_b(x),\rho_b(x)\mid b\bigr)+O(1).
\]
再由链式不等式（吸收自描述开销）存在常数 \(C_1\) 使
\[
K\bigl(\pi_b(x),\rho_b(x)\mid b\bigr)
\le
K\bigl(\pi_b(x)\mid b\bigr)+K\bigl(\rho_b(x)\mid \pi_b(x),b\bigr)+C_1\log b.
\]
又因 \(\pi_b(x)\in \mathrm{Phase}(k)_b\) 且 \(|\mathrm{Phase}(k)_b|=2^{bk}\)，存在常数 \(C_2\) 使
\[
K\bigl(\pi_b(x)\mid b\bigr)\le bk+C_2.
\]
合并并整理，得到
\[
K\bigl(\rho_b(x)\mid \pi_b(x),b\bigr)
\ge
(r-k)b+\log_2|T|-t-C\log b,
\]
其中 \(C\) 吸收 \(C_0,C_1,C_2\) 与 \(O(1)\) 项。
因此满足严格小于该下界者只能落在上述例外集合中，其大小被 \(2^{-t}N_b\) 控制，计数与概率式同时成立。证毕。
\end{proof}

\begin{corollary}[单相位通道的指数残差下界：缺失圆因子以条件算法信息外置]\label{cor:xi-cdim-kolm-single-phase}
在定理 \ref{thm:xi-cdim-kolm-residual-gap} 的设定下取 \(k=1\) 且 \(r\ge 2\)。
则对任意 \(t\ge 1\)，除至多 \(2^{-t}\) 比例的窗口元素外，均满足
\[
K\bigl(\rho_b(x)\mid \pi_b(x),b\bigr)\ \ge\ (r-1)b+\log_2|T|-t-C\log b.
\]
\end{corollary}
\begin{proof}
由定理 \ref{thm:xi-cdim-kolm-residual-gap} 直接代入 \(k=1\) 即得。证毕。
\end{proof}

\begin{example}[Basic examples / 基本例子]
\label{ex:cdim-signature}
\begin{enumerate}
  \item \textbf{整数与 \(k\) 维整数格.}
        由例 \ref{ex:circle-dimension-examples} 有 \(\cdim(\ZZ)=1\)、\(\cdim(\ZZ^k)=k\)。
        纯可逆整数寄存器在无残差时可取 \(\Sigma_{\mathrm{res}}=\varnothing\)，故
        \[
        \Sigma\Sigma(\ZZ)=\bigl((1,0),\varnothing\bigr).
        \]
  \item \textbf{自然数的“半圆维”.}
        把 \(M=\NN\) 视为 \(\ZZ\) 的正锥，则 \(G(M)=\ZZ\)，
        从而 \(\cdim_{+}(\NN)=\tfrac12\)（例 \ref{ex:circle-dimension-examples}）。
  \item \textbf{有限标签（纤维/离散细分）.}
        对任意有限交换群 \(F\)，有 \(\cdim(F)=0\)（例 \ref{ex:circle-dimension-examples}）。
        因此有限标签不会出现在 \(\Sigma_{\circ}\) 中，而应进入 \(\Sigma_{\mathrm{res}}\)。
  \item \textbf{有限纤维扩张不增圆维.}
        若存在短正合列 \(0\to F\to G\to H\to 0\) 且 \(F\) 有限，
        则 \(\cdim(G)=\cdim(H)\)（命题 \ref{prop:circle-dimension-laws}）。
        于是“同一圆维、不同标签”的差异被 \(\Sigma_{\mathrm{res}}\) 捕获，而非由圆维区分。
  \item \textbf{相位门的残差寄存器模板.}
        当实现引入截断/量化/折叠等不可逆环节时，可按
        \(
        \Sigma_{\mathrm{res}}(X)=(\mathsf{End}_\varepsilon,\mathsf{Prec},\mathsf{Tag})
        \)
        记录端点窗口质量、精度势累计与不可寻址通道标记（\S\ref{subsec:cdim-implementation-audit}），从而维持账本闭合。
\end{enumerate}
\end{example}

\paragraph{圆维谱：有限 torsion 账本的相位—寄存器 Pareto 前沿}
本段在定义 \ref{def:circle-dimension}--\ref{def:half-circle-dimension} 的圆维记号背景下，给出一类完全可计算的“相位圆数 \(\leftrightarrow\) 离散账本规模”权衡谱。
其目标并非改变结构圆维（有限标签在结构层恒满足 \(\cdim=0\)），而是把实现层的“可逆分流”约束压缩为一条可逆恢复的资源谱，从而将有限交换账本的同构类型以最小剩余函数的形式显式化。

\begin{definition}[相位—寄存器实现与最小剩余函数]\label{def:xi-cdim-minimal-residual-function}
设 \(F\) 为有限交换群，\(k\in\NN\cup\{0\}\)。
称 \((\alpha,\beta)\) 为 \(F\) 的一个\textbf{\(k\)-相位—寄存器实现}，若存在有限交换群 \(R\) 及群同态
\[
\alpha:F\to \TT^{k},\qquad \beta:F\to R
\]
使得联合映射
\[
(\alpha,\beta):F\to \TT^{k}\times R
\]
为单射。
\par
定义 \(F\) 的\textbf{最小剩余寄存器阶函数}为
\[
\Lambda_F(k):=\min\Bigl\{|R|:\ \exists\,(\alpha,\beta)\ \text{为 }k\text{-相位—寄存器实现}\Bigr\}\in\NN.
\]
并定义 \(F\) 的\textbf{圆维谱}为有限序列
\[
\operatorname{Spec}_{\circ}(F):=\bigl(\Lambda_F(0),\Lambda_F(1),\dots,\Lambda_F(t)\bigr),
\qquad
t:=\min\{k\ge 0:\Lambda_F(k)=1\}.
\]
\end{definition}

\begin{lemma}[有限交换群的 \(k\)-生成商最大阶]\label{lem:xi-cdim-k-generated-quotient-max-order}
\leanverified{paper\_xi\_cdim\_k\_generated\_quotient\_max\_order}
设 \(F\cong \ZZ_{n_1}\oplus\cdots\oplus\ZZ_{n_t}\) 为有限交换群的不变量因子分解（\(1<n_1\mid\cdots\mid n_t\)）。
若 \(Q\) 为 \(F\) 的商群且可由不超过 \(k\) 个元素生成，则
\[
|Q|\ \le\ \prod_{i=t-k+1}^{t}n_i.
\]
并且该上界由投影商 \(F\twoheadrightarrow \ZZ_{n_{t-k+1}}\oplus\cdots\oplus\ZZ_{n_t}\) 达到。
\end{lemma}
\begin{proof}
达到性由尾部投影立即给出，故只需证明上界。
\par
设 \(Q\cong F/K\)。对任意素数 \(p\) 记 \(F_p\)、\(Q_p\) 为其 \(p\)-primary 分量，则 \(Q_p\) 为 \(F_p\) 的商群且同样可由不超过 \(k\) 个元素生成。
写
\(
e_{i,p}:=v_p(n_i)\in\ZZ_{\ge 0}
\)。
则
\[
F_p\ \cong\ \ZZ_{p^{e_{1,p}}}\oplus\cdots\oplus\ZZ_{p^{e_{t,p}}},\qquad
0\le e_{1,p}\le\cdots\le e_{t,p}.
\]
对有限 \(p\)-群 \(A\) 令
\[
a_r(A):=\dim_{\FF_p}\bigl(p^{r-1}A/p^rA\bigr)\qquad(r\ge 1).
\]
由链 \(A\supset pA\supset p^2A\supset\cdots\supset 0\)，有
\(
|A|=\prod_{r\ge 1}\bigl|p^{r-1}A/p^rA\bigr|
\)
从而
\begin{equation}\label{eq:xi-cdim-kquot-logp-sum}
\log_p|A|=\sum_{r\ge 1}a_r(A).
\end{equation}
对 \(A=F_p\)，由直和分解可知
\(
a_r(F_p)=\#\{i:\ e_{i,p}\ge r\}
\)。
对 \(A=Q_p\)，商映射 \(F_p\twoheadrightarrow Q_p\) 在每个层级诱导满射
\(
p^{r-1}F_p/p^rF_p \twoheadrightarrow p^{r-1}Q_p/p^rQ_p
\)，
故
\(
a_r(Q_p)\le a_r(F_p)
\)。
又因 \(Q_p\) 可由 \(\le k\) 个元素生成，故
\(
a_1(Q_p)=\dim_{\FF_p}(Q_p/pQ_p)\le k
\)，并由单调性 \(a_{r+1}(Q_p)\le a_r(Q_p)\) 得对所有 \(r\ge 1\) 有 \(a_r(Q_p)\le k\)。
于是结合 \eqref{eq:xi-cdim-kquot-logp-sum} 得
\[
\log_p|Q_p|
\le
\sum_{r\ge 1}\min\{a_r(F_p),k\}.
\]
最后记 \(a_r:=a_r(F_p)=\#\{i:\ e_{i,p}\ge r\}\)。
有恒等式
\begin{equation}\label{eq:xi-cdim-min-sum-largest-exponents}
\sum_{r\ge 1}\min\{a_r,k\}=\sum_{i=t-k+1}^{t}e_{i,p},
\end{equation}
其证明为：将左端改写为
\(
\sum_{r\ge 1}\sum_{j=1}^{k}\ind_{\{a_r\ge j\}}
=\sum_{j=1}^{k}\sum_{r\ge 1}\ind_{\{a_r\ge j\}}
\)；
而 \(a_r\ge j\) 当且仅当第 \(t-j+1\) 个最大指数满足 \(e_{t-j+1,p}\ge r\)，故内层求和等于 \(e_{t-j+1,p}\)，从而得到 \eqref{eq:xi-cdim-min-sum-largest-exponents}。
因此
\(
|Q_p|\le p^{\sum_{i=t-k+1}^{t}e_{i,p}}=\prod_{i=t-k+1}^{t}p^{e_{i,p}}
\)。
对所有素数相乘即得
\[
|Q|=\prod_{p}|Q_p|
\le
\prod_{p}\prod_{i=t-k+1}^{t}p^{e_{i,p}}
\;=\;
\prod_{i=t-k+1}^{t}\prod_{p}p^{v_p(n_i)}
\;=\;
\prod_{i=t-k+1}^{t}n_i,
\]
完成证明。
\end{proof}

\begin{theorem}[闭式 Pareto 前沿：\texorpdfstring{$\Lambda_F(k)$}{Lambda(k)} 的显式公式]\label{thm:xi-cdim-lambda-closed-form}
设有限交换群 \(F\) 的不变量因子分解为
\[
F\ \cong\ \ZZ_{n_1}\oplus\cdots\oplus\ZZ_{n_t},
\qquad
1<n_1\mid n_2\mid\cdots\mid n_t,
\]
其中 \((n_1,\dots,n_t)\) 唯一确定。
则对任意 \(k\ge 0\)，有
\[
\Lambda_F(k)=
\begin{cases}
\displaystyle \prod_{i=1}^{t-k} n_i,& 0\le k\le t,\\[6pt]
1,& k\ge t.
\end{cases}
\]
特别地，\(\Lambda_F(k)=1\) 当且仅当 \(k\ge t\)，从而定义 \ref{def:xi-cdim-minimal-residual-function} 中的 \(t\) 等于 \(F\) 的不变量因子个数（亦即最小生成元数）。
\leanverified{paper\_xi\_cdim\_lambda\_closed\_form}
\end{theorem}
\begin{proof}
取任意 \(k\)-相位—寄存器实现 \((\alpha,\beta)\) 及其残余群 \(R\)，并记
\(
K:=\ker(\alpha)\le F
\)。
由于 \((\alpha,\beta)\) 单射，\(\beta|_{K}:K\to R\) 亦为单射，故
\begin{equation}\label{eq:xi-cdim-lambda-lb1}
|R|\ \ge\ |K|.
\end{equation}
另一方面，\(\alpha(F)\) 是 \(\TT^{k}\) 的有限子群，因而同构于某个 \((\mu_N)^k\cong(\ZZ_{N})^{k}\) 的子群（对某个 \(N\)），从而可由不超过 \(k\) 个元素生成。
记 \(Q:=\alpha(F)\cong F/K\)。
由引理 \ref{lem:xi-cdim-k-generated-quotient-max-order}（注意 \(Q\) 可由不超过 \(k\) 个元素生成）得
\(
|Q|\le \prod_{i=t-k+1}^{t}n_i
\)。
由 \(|K|=|F|/|Q|\) 与 \eqref{eq:xi-cdim-lambda-lb1} 得
\[
|R|\ \ge\ |K|
\ \ge\
\frac{\prod_{i=1}^{t}n_i}{\prod_{i=t-k+1}^{t}n_i}
\ =\
\prod_{i=1}^{t-k}n_i.
\]
从而 \(\Lambda_F(k)\ge \prod_{i=1}^{t-k}n_i\)（当 \(k\le t\)）。
\par
为构造达到等号的实现，写
\[
F\ \cong\ K_0\oplus Q_0,\qquad
K_0:=\ZZ_{n_1}\oplus\cdots\oplus\ZZ_{n_{t-k}},\quad
Q_0:=\ZZ_{n_{t-k+1}}\oplus\cdots\oplus\ZZ_{n_t}.
\]
令 \(\alpha\) 为投影到 \(Q_0\) 后将每个循环因子嵌入一维圆的根单位子群（从而得到 \(Q_0\hookrightarrow \TT^{k}\)）；
令 \(R:=K_0\) 且 \(\beta\) 为投影到 \(K_0\)。
则 \((\alpha,\beta)\) 显然为单射，且 \(|R|=|K_0|=\prod_{i=1}^{t-k}n_i\)，从而下界可达，给出所示闭式。
当 \(k\ge t\) 时取 \(R\) 为平凡群即得 \(\Lambda_F(k)=1\)。
\end{proof}

\begin{corollary}[圆维谱的完备性：由 \texorpdfstring{$\operatorname{Spec}_{\circ}(F)$}{Spec(F)} 恢复不变量因子]\label{cor:xi-cdim-spectrum-completeness}
\leanverified{paper\_xi\_cdim\_spectrum\_completeness}
沿用定理 \ref{thm:xi-cdim-lambda-closed-form} 的记号。
令 \(t=\min\{k:\Lambda_F(k)=1\}\)。
则对 \(k=0,1,\dots,t-1\) 有递推恒等式
\[
n_{t-k}=\frac{\Lambda_F(k)}{\Lambda_F(k+1)}.
\]
因此 \(\operatorname{Spec}_{\circ}(F)\) 与不变量因子序列 \((n_1,\dots,n_t)\) 互相可逆确定，从而决定 \(F\) 的同构类型。
\end{corollary}
\begin{proof}
由定理 \ref{thm:xi-cdim-lambda-closed-form}，
\(
\Lambda_F(k)=\prod_{i=1}^{t-k}n_i
\)
与
\(
\Lambda_F(k+1)=\prod_{i=1}^{t-k-1}n_i
\)
相除即得。
其余断言由不变量因子分解唯一性立得。
\end{proof}

\begin{proposition}[有限生成交换群的圆维—账本可计算签名]\label{prop:xi-cdim-circle-ledger-signature}
\leanverified{paper\_xi\_cdim\_circle\_ledger\_signature}
对任意有限生成交换群 \(G\)，记其 torsion 子群为
\(
G_{\mathrm{tor}}:=\{x\in G:\ \exists\,n\ge 1,\ nx=0\}
\)。
定义
\[
\Sigma_{\circ}^{\mathrm{ab}}(G)
:=\Bigl(\cdim(G),\ \operatorname{Spec}_{\circ}(G_{\mathrm{tor}})\Bigr).
\]
则：
\begin{enumerate}
  \item \(\Sigma_{\circ}^{\mathrm{ab}}(G)\) 可由 \(G\cong \ZZ^{d}\oplus G_{\mathrm{tor}}\) 的分解计算得到，并满足 \(\cdim(G)=d\)；
  \item \(\Sigma_{\circ}^{\mathrm{ab}}(G)\) 决定 \(G\) 的同构类型：若 \(G,H\) 为有限生成交换群，则
  \[
  \Sigma_{\circ}^{\mathrm{ab}}(G)=\Sigma_{\circ}^{\mathrm{ab}}(H)
  \quad\Longleftrightarrow\quad
  G\cong H.
  \]
\end{enumerate}
\end{proposition}
\begin{proof}
有限生成交换群结构定理给出 \(G\cong \ZZ^{d}\oplus G_{\mathrm{tor}}\)，且该分解在同构意义下唯一。
由定义 \ref{def:circle-dimension}，\(\widehat{\ZZ^{d}}\cong \TT^{d}\) 的连通分量维数为 \(d\)，而 \(\widehat{G_{\mathrm{tor}}}\) 有限从而不贡献连通分量维数，故 \(\cdim(G)=d\)。
再由推论 \ref{cor:xi-cdim-spectrum-completeness}，\(\operatorname{Spec}_{\circ}(G_{\mathrm{tor}})\) 可逆确定 \(G_{\mathrm{tor}}\) 的不变量因子，从而确定 torsion 部分的同构类型；与秩 \(d\) 合并即得结论。
\end{proof}

\begin{definition}[圆维谱配分多项式]\label{def:xi-cdim-spectrum-partition-polynomial}
对有限交换群 \(F\)（记 \(t=\min\{k:\Lambda_F(k)=1\}\)），定义其圆维谱配分多项式为
\[
Z_F(z):=\sum_{k=0}^{t}\Lambda_F(k)\,z^{k}\ \in\ \RR[z].
\]
\end{definition}

\begin{theorem}[零点环带界：圆维谱配分多项式的 Lee--Yang 型几何定律]\label{thm:xi-cdim-spectrum-zeros-annulus}
设 \(F\cong \ZZ_{n_1}\oplus\cdots\oplus\ZZ_{n_t}\) 为定理 \ref{thm:xi-cdim-lambda-closed-form} 的不变量因子分解，并令 \(Z_F\) 如定义 \ref{def:xi-cdim-spectrum-partition-polynomial}。
则 \(Z_F\) 的全部零点满足环带约束
\[
n_1\ \le\ |z|\ \le\ n_t.
\]
\leanverified{paper\_xi\_cdim\_spectrum\_zeros\_annulus}
\end{theorem}
\begin{proof}
令 \(a_k:=\Lambda_F(k)\)（\(0\le k\le t\)）。由定理 \ref{thm:xi-cdim-lambda-closed-form} 得
\[
\frac{a_k}{a_{k+1}}=\frac{\Lambda_F(k)}{\Lambda_F(k+1)}=n_{t-k}\qquad(0\le k<t),
\]
从而
\[
\min_{0\le k<t}\frac{a_k}{a_{k+1}}=n_1,\qquad
\max_{0\le k<t}\frac{a_k}{a_{k+1}}=n_t.
\]
对具有正系数的多项式 \(P(z)=\sum_{k=0}^{t}a_k z^k\)，经典 Enestr\"om--Kakeya 环带定理断言：其零点均落在
\(
\min_{k}\frac{a_k}{a_{k+1}}\le |z|\le \max_{k}\frac{a_k}{a_{k+1}}
\)
所界定的环带内。
将 \(P=Z_F\) 代入即得结论。
\end{proof}

\begin{corollary}[同阶均匀 torsion 的圆定律]\label{cor:xi-cdim-spectrum-homocyclic-circle-law}
若 \(F\cong (\ZZ_{n})^{m}\)（\(n>1\)），则
\[
Z_F(z)=\sum_{k=0}^{m} n^{m-k}z^{k}=n^{m}\sum_{k=0}^{m}(z/n)^{k},
\]
其零点恰为
\[
z_j=n\,e^{\frac{2\pi i j}{m+1}},\qquad j=1,2,\dots,m.
\]
特别地，当 \(n=p\) 为素数且 \(F\cong (\ZZ_{p})^{m}\) 时，零点严格落在圆 \(|z|=p\) 上并等角分布。
\leanverified{paper\_xi\_cdim\_spectrum\_homocyclic\_circle\_law}
\end{corollary}
\begin{proof}
由几何级数闭式 \(1+w+\cdots+w^m=(w^{m+1}-1)/(w-1)\) 得
\(
Z_F(z)=n^{m}\frac{(z/n)^{m+1}-1}{z/n-1}
\)。
因此零点由 \((z/n)^{m+1}=1\) 且 \(z\ne n\) 给出，即所示集合。
\end{proof}

\begin{example}[典型对象的“同圆维不同谱”判别]\label{ex:xi-cdim-same-cdim-different-spectrum}
\begin{enumerate}
  \item \textbf{有限域的加法群.}
        对任意素数 \(p\) 与 \(m\ge 1\)，有 \((\FF_{p^m},+)\cong (\ZZ_{p})^{m}\)。
        因而由推论 \ref{cor:xi-cdim-spectrum-homocyclic-circle-law}，其 \(Z\)-零点严格落在 \(|z|=p\) 上，并由半径 \(p\) 给出一个仅依赖底素数的刚性几何指纹。
  \item \textbf{类群与任意有限 torsion 账本.}
        对任意数域 \(K\)，理想类群 \(\mathrm{Cl}(K)\) 为有限交换群且 \(\cdim(\mathrm{Cl}(K))=0\)。
        推论 \ref{cor:xi-cdim-spectrum-completeness} 表明：即便结构圆维同为零，\(\operatorname{Spec}_{\circ}(\mathrm{Cl}(K))\) 仍可作为同构类型的完备离散签名。
  \item \textbf{同圆维的严格区分范式.}
        令 \(G_1=\ZZ^{d}\oplus \ZZ_{n}\)、\(G_2=\ZZ^{d}\oplus \ZZ_{m}\)。
        则 \(\cdim(G_1)=\cdim(G_2)=d\)，但
        \(
        \operatorname{Spec}_{\circ}(\ZZ_{n})=(n,1)\neq(m,1)=\operatorname{Spec}_{\circ}(\ZZ_{m})
        \)
        当 \(n\neq m\)，从而 \(\Sigma_{\circ}^{\mathrm{ab}}(G_1)\neq \Sigma_{\circ}^{\mathrm{ab}}(G_2)\)。
\end{enumerate}
\end{example}

\paragraph{圆维的可观测重构：素数局部账本、解析极点与相位精度指数}
以下结果把结构层圆维 \(\cdim\) 与实现层“有限分辨率可逆性”联结到同一组可计算对象上：
一方面，\(|G/nG|\) 形成 Dirichlet 解析账本并将 \(\cdim\) 编码为收敛临界指数；
另一方面，\(\ZZ^r\to\TT^d\) 的相位压缩必然产生 \(B^{-r/d}\) 量级的最小分离衰减，从而给出不可规避的精度守恒律。

\begin{theorem}[圆维 \texorpdfstring{$\zeta$}{zeta} 函数的欧拉分解与收敛界]\label{thm:xi-cdim-zeta-euler-abscissa-locality}
\leanverified{paper\_xi\_cdim\_zeta\_euler\_abscissa\_locality}
设 \(G\) 为有限生成离散交换群，并记
\[
r:=\cdim(G)\in\NN.
\]
对每个 \(n\ge 1\) 定义
\[
a_G(n):=\card{G/nG},
\qquad
\zeta_{\circ}(G;s):=\sum_{n\ge 1}a_G(n)\,n^{-s}.
\]
则成立：
\begin{enumerate}
  \item \(a_G\) 为乘法函数，故在收敛半平面内
  \[
  \zeta_{\circ}(G;s)=\prod_{p}\left(\sum_{k\ge 0}a_G(p^k)\,p^{-ks}\right).
  \]
  \item \(\zeta_{\circ}(G;\cdot)\) 的收敛 abscissa 为
  \[
  \sigma_c\!\bigl(\zeta_{\circ}(G;\cdot)\bigr)=r+1.
  \]
  \item 若 \(G\cong \ZZ^r\oplus F\)（\(F\) 有限），则存在仅依赖于有限坏素数集
  \[
  S(G):=\{p:\ p\mid |F|\}
  \]
  的欧拉修正因子 \(E_G(s)\)，使
  \[
  \zeta_{\circ}(G;s)=\zeta(s-r)\,E_G(s),
  \]
  其中 \(E_G(s)\) 为 \(S(G)\) 上有理局部因子的有限乘积。
\end{enumerate}
\end{theorem}
\begin{proof}
先证乘法性。取 \(\gcd(m,n)=1\)，考虑
\[
\phi:G\to G/mG\times G/nG,\qquad g\mapsto (g\bmod mG,\ g\bmod nG).
\]
有 \(\ker\phi=mG\cap nG\)。由 Bezout 等式取 \(u,v\in\ZZ\) 使 \(um+vn=1\)。
若 \(g\in mG\cap nG\)，写 \(g=ma=nb\)，则
\[
g=(um+vn)g=um(nb)+vn(ma)=mn(ub+va)\in mnG,
\]
故 \(mG\cap nG=mnG\)。
另一方面，任取 \((x\bmod mG,\ y\bmod nG)\)，令 \(g:=vnx+umy\)，则
\[
g-x=um(y-x)\in mG,\qquad g-y=vn(x-y)\in nG,
\]
故 \(\phi\) 满射。
于是 \(G/mnG\cong G/mG\times G/nG\)，即
\[
a_G(mn)=a_G(m)a_G(n),
\]
故 \(a_G\) 为乘法函数，欧拉积成立。

再证收敛界。写 \(G\cong \ZZ^r\oplus F\)（\(F\) 有限），则
\[
G/nG\cong (\ZZ/n\ZZ)^r\oplus (F/nF),
\qquad
a_G(n)=n^r\,|F/nF|.
\]
由 \(1\le |F/nF|\le |F|\)，得
\[
\sum_{n\ge 1}n^{r-s}
\le
\zeta_{\circ}(G;s)
\le
|F|\sum_{n\ge 1}n^{r-s}.
\]
故其收敛当且仅当 \(\Re(s)>r+1\)，从而 \(\sigma_c=r+1\)。

最后证局部性。对 \(p\nmid |F|\)，乘以 \(p^k\) 在 \(F\) 上为自同构，故 \(F/p^kF\) 平凡，因而
\[
\sum_{k\ge 0}a_G(p^k)p^{-ks}
=
\sum_{k\ge 0}p^{k(r-s)}
=
\frac{1}{1-p^{r-s}}.
\]
把这些好素数局部因子相乘得到 \(\zeta(s-r)\)；坏素数仅在 \(S(G)\) 中有限多个，其局部比值合并为有限欧拉修正 \(E_G(s)\)。证毕。
\end{proof}

\begin{remark}[同态计数口径与 \texorpdfstring{$n$}{n}-层商口径的一致性]\label{rem:xi-cdim-hom-quotient-consistency}
在有限生成离散交换群的范围内，循环靶上的同态计数与模 \(n\) 商的基数严格一致。
具体地，对任意有限生成离散交换群 \(G\) 与任意 \(n\ge 1\)，有自然双射
\[
\operatorname{Hom}(G,\ZZ/n\ZZ)\ \cong\ \operatorname{Hom}(G,\mu_n),
\]
且由命题 \ref{prop:cdim-phase-spectrum-quotient}，
\[
\card{\operatorname{Hom}(G,\ZZ/n\ZZ)}=\card{\operatorname{Hom}(G,\mu_n)}=\card{G/nG}.
\]
因此，若以
\(
a_n(G):=\card{\operatorname{Hom}(G,\ZZ/n\ZZ)}
\)
定义 Dirichlet 级数
\(
\zeta_G(s):=\sum_{n\ge 1}a_n(G)n^{-s}
\)，
则 \(\zeta_G=\zeta_{\circ}(G;\cdot)\) 与定理 \ref{thm:xi-cdim-zeta-euler-abscissa-locality} 的欧拉分解与收敛界完全一致。
\end{remark}

\begin{corollary}[圆维 \texorpdfstring{$\zeta$}{zeta} 指纹的同构刚性]\label{cor:xi-cdim-zeta-fingerprint-rigidity}
\leanverified{paper\_xi\_cdim\_zeta\_fingerprint\_rigidity}
设 \(G,H\) 为有限生成离散交换群。
若对一切 \(n\ge 1\) 都有
\[
\card{\operatorname{Hom}(G,\ZZ/n\ZZ)}=\card{\operatorname{Hom}(H,\ZZ/n\ZZ)},
\]
则 \(G\cong H\)。
等价地，若 \(\zeta_G(s)\equiv \zeta_H(s)\) 作为 Dirichlet 级数一致，则 \(G\cong H\)。
\end{corollary}
\begin{proof}
由注 \ref{rem:xi-cdim-hom-quotient-consistency}，
\(
\card{\operatorname{Hom}(G,\ZZ/n\ZZ)}=\card{G/nG}
\)
与
\(
\card{\operatorname{Hom}(H,\ZZ/n\ZZ)}=\card{H/nH}
\)。
故假设等价于 \(|G/nG|=|H/nH|\) 对一切 \(n\ge 1\) 成立。
由定理 \ref{thm:xi-cdim-prime-ledger-complete-reconstruction} 立即推出 \(G\cong H\)。
\end{proof}

\begin{proposition}[素数局部反演：离散导数给出 \texorpdfstring{$p$}{p}-primary 不变量因子]\label{prop:xi-cdim-local-inversion-delta}
设 \(G\cong \ZZ^{r}\oplus F\) 为有限生成离散交换群，固定素数 \(p\)。
对 \(k\ge 0\) 定义
\[
c_p(k):=v_p\!\bigl(\card{G/p^kG}\bigr)-kr,\qquad c_p(0):=0,
\]
并令离散导数 \(\Delta_p(k):=c_p(k)-c_p(k-1)\)（\(k\ge 1\)）。
将 \(F_p\) 写成
\[
F_p\cong \bigoplus_{j=1}^{t_p}\ZZ/p^{\lambda_{p,j}}\ZZ,\qquad
\lambda_{p,1}\ge\cdots\ge \lambda_{p,t_p}\ge 1.
\]
则：
\begin{enumerate}
  \item \(\Delta_p(k)\in\NN\) 且随 \(k\) 单调不增；
  \item 当 \(k>\max_j\lambda_{p,j}\) 时 \(\Delta_p(k)=0\)；
  \item 不变因子指数多重集由反演公式唯一给出：
  \[
  \#\{j:\lambda_{p,j}=k\}=\Delta_p(k)-\Delta_p(k+1)\qquad(k\ge 1).
  \]
\end{enumerate}
\end{proposition}
\begin{proof}
定理 \ref{thm:xi-cdim-prime-ledger-complete-reconstruction} 的证明已经给出恒等式
\[
c_p(k)=v_p(|F_p/p^kF_p|)=\sum_{j=1}^{t_p}\min(\lambda_{p,j},k).
\]
于是
\[
\Delta_p(k)=\#\{j:\lambda_{p,j}\ge k\},
\]
从而单调性与终局为零成立。
最后一式为”列高度函数反演”：\(\#\{\lambda_{p,j}=k\}=\#\{\lambda_{p,j}\ge k\}-\#\{\lambda_{p,j}\ge k+1\}\)。
\leanverified{filter\_ge\_card\_mono}
\leanverified{filter\_ge\_card\_sub\_succ}
\leanverified{paper\_xi\_cdim\_local\_inversion\_delta}
\end{proof}

\begin{proposition}[有限局部数据的充分性]\label{prop:xi-cdim-finite-data-suffices}
\leanverified{paper\_xi\_cdim\_finite\_data\_suffices}
设 \(G\cong \ZZ^{r}\oplus F\) 为有限生成离散交换群，并记 \(e:=\exp(F)\)。
则 \(G\) 的同构类型由以下有限数据完全确定：
\begin{itemize}
  \item 自由秩 \(r=\cdim(G)\)；
  \item 对每个素数 \(p\mid e\)，有限序列 \(\{\card{G/p^kG}\}_{1\le k\le v_p(e)}\)。
\end{itemize}
并且对任意 \(p\nmid e\) 与任意 \(k\ge 1\) 恒有 \(\card{G/p^kG}=p^{kr}\)。
\end{proposition}
\begin{proof}
当 \(p\nmid e\) 时，乘以 \(p^k\) 在 \(F\) 上为自同构，故 \(F/p^kF\) 平凡，从而
\(\card{G/p^kG}=p^{kr}\)（亦见命题 \ref{prop:xi-cdim-good-prime-readout-and-bad-support}）。
当 \(p\mid e\) 时，\(F_p\) 的最大指数不超过 \(v_p(e)\)；由命题 \ref{prop:xi-cdim-local-inversion-delta}，
\(\{\card{G/p^kG}\}_{1\le k\le v_p(e)}\) 已足以反演得到 \(F_p\) 的不变因子。
汇总所有 \(p\) 的 \(p\)-primary 分量并与自由秩 \(r\) 合并，即唯一确定 \(G\) 的同构类型。
\leanpartial{paper\_xi\_cdim\_finite\_quotient\_free\_part}{仅自由部分 $|(\ZZ/n\ZZ)^r| = n^r$；完整 fg abelian group 结构定理反演待后轮}
\leanverified{card\_zmod}
\leanverified{card\_zmod\_pi}
\end{proof}

\begin{theorem}[有限素数集局部化的 \texorpdfstring{$\zeta$}{zeta} 公式]\label{thm:xi-cdim-localization-zeta-AS}
\leanverified{paper\_xi\_cdim\_localization\_zeta\_as}
取有限素数集 \(S\)，令 \(A_S:=\ZZ[S^{-1}]\) 视为加法群。
定义
\[
a_{A_S}(n):=\card{\operatorname{Hom}(A_S,\ZZ/n\ZZ)},\qquad
\zeta_{A_S}(s):=\sum_{n\ge 1}a_{A_S}(n)\,n^{-s}.
\]
则对所有 \(n\ge 1\) 有
\[
a_{A_S}(n)=n_{S^c}:=\prod_{\substack{p^k\parallel n\\ p\notin S}}p^k,
\]
从而 \(\zeta_{A_S}\) 具有欧拉分解并满足恒等式
\[
\zeta_{A_S}(s)
=\prod_{p\notin S}\frac{1}{1-p^{1-s}}\cdot\prod_{p\in S}\frac{1}{1-p^{-s}}
=\zeta(s-1)\cdot\prod_{p\in S}\frac{1-p^{1-s}}{1-p^{-s}}.
\]
\end{theorem}
\begin{proof}
任意群同态 \(f:A_S\to \ZZ/n\ZZ\) 由 \(1\in A_S\) 的像唯一决定，且必须满足：对每个 \(p\in S\)，元素 \(p^{-1}\in A_S\) 在 \(\ZZ/n\ZZ\) 中的像必须存在。
因此当且仅当 \(p\) 在 \(\ZZ/n\ZZ\) 中可逆时（即 \(p\nmid n\)），该素数才对可选像施加约束；若 \(p\mid n\)，则该局部因子在计数上仅贡献常数项。
更直接地，作为环（因而也作为加法群）有同构
\[
A_S/nA_S\cong (\ZZ/n\ZZ)[S^{-1}],
\]
并利用中国剩余分解
\(
\ZZ/n\ZZ\cong \prod_{p^k\parallel n}\ZZ/p^k\ZZ
\)
可得
\[
(\ZZ/n\ZZ)[S^{-1}]\cong \prod_{\substack{p^k\parallel n\\ p\notin S}}\ZZ/p^k\ZZ,
\]
故
\(
\card{A_S/nA_S}=n_{S^c}
\)。
由注 \ref{rem:xi-cdim-hom-quotient-consistency} 的同态计数一致性，得到 \(a_{A_S}(n)=\card{A_S/nA_S}\)，从而 \(a_{A_S}(n)=n_{S^c}\)。
欧拉分解由 \(n_{S^c}\) 的完全乘法性给出；对素幂 \(n=p^k\) 的两种情形分别求和得到局部因子
\[
\sum_{k\ge 0}\frac{a_{A_S}(p^k)}{p^{ks}}
=
\begin{cases}
\sum_{k\ge 0}\frac{p^k}{p^{ks}}=\frac{1}{1-p^{1-s}},& p\notin S,\\[6pt]
\sum_{k\ge 0}\frac{1}{p^{ks}}=\frac{1}{1-p^{-s}},& p\in S.
\end{cases}
\]
乘积即得所示恒等式。
\end{proof}

\begin{theorem}[分母轴的直接极限与一维溶子连通因子]\label{thm:xi-denominator-axis-direct-limit-solenoid}
\leanverified{paper\_xi\_denominator\_axis\_direct\_limit\_solenoid}
给定整数序列 \((d_n)_{n\ge 0}\subset\ZZ_{\ge 2}\)，并定义离散交换群
\[
A:=\mathrm{colim}\bigl(\ZZ \xrightarrow{\times d_0}\ZZ \xrightarrow{\times d_1}\ZZ \xrightarrow{\times d_2}\cdots\bigr).
\]
令
\[
\mathcal P:=\{p\ \text{素数}:\ \exists n,\ p\mid d_n\},
\qquad
A_{\mathcal P}:=\ZZ[\mathcal P^{-1}]\subset\QQ
\]
为对 \(\mathcal P\) 的局部化（按加法群理解）。
则存在自然同构 \(A\cong A_{\mathcal P}\)。
此外，Pontryagin 对偶 \(\widehat A\) 可识别为圆群在自同态 \(z\mapsto z^{d_n}\) 下的逆极限：
\[
\widehat A\ \cong\ \varprojlim\bigl(\TT \xleftarrow{(\cdot)^{d_0}}\TT \xleftarrow{(\cdot)^{d_1}}\TT \xleftarrow{(\cdot)^{d_2}}\cdots\bigr),
\]
其单位连通分量 \((\widehat A)^0\) 为经典一维 solenoid；因此 \((\widehat A)^0\) 为紧、连通、交换拓扑群，且覆盖维数为 \(1\)。
\end{theorem}
\begin{proof}
按定义，\(A\) 由同态 \(\ZZ\to\QQ\)（把第 \(n\) 层的 \(1\) 送到 \(1/D_n\)，其中 \(D_n:=\prod_{j<n}d_j\)）给出一个自然嵌入到 \(\QQ\) 的模型，从而 \(A\) 由形如 \(m/D_n\) 的有理数生成。
其分母的素因子恰落在集合 \(\mathcal P\) 内，故 \(A\subseteq \ZZ[\mathcal P^{-1}]\)。
反过来，对任意 \(p\in\mathcal P\) 存在 \(n\) 使 \(p\mid d_n\)，从而 \(1/p\in A\)；
因此 \(A\) 包含 \(\mathcal P\) 中全部素数的倒数并对整数封闭，推出 \(\ZZ[\mathcal P^{-1}]\subseteq A\)，故 \(A\cong \ZZ[\mathcal P^{-1}]\)。
\par
对偶描述来自一般事实：离散群的直接极限对偶为紧群的逆极限。
并且 \(\ZZ[\mathcal P^{-1}]\) 的对偶是由圆群在幂映射下的逆极限构成的 solenoid，其单位连通分量为一维（覆盖维数 \(1\)）的连通紧阿贝尔群。证毕。
\end{proof}

\begin{proposition}[素数支持的内禀刻画：满射性与模消失]\label{prop:xi-denominator-axis-prime-support-characterization}
\leanverified{paper\_xi\_denominator\_axis\_prime\_support\_characterization}
沿用定理 \ref{thm:xi-denominator-axis-direct-limit-solenoid} 的记号。
对任意素数 \(p\)，以下命题等价：
\begin{enumerate}
  \item \(p\in\mathcal P\)；
  \item \(A=pA\)（即乘以 \(p\) 的自同态 \(A\to A\) 为满射）；
  \item \(A/pA=0\)。
\end{enumerate}
因此 \(\mathcal P\) 可由抽象阿贝尔群 \(A\) 内禀恢复。
\end{proposition}
\begin{proof}
(1)\(\Rightarrow\)(2)：若 \(p\in\mathcal P\)，则 \(1/p\in A\)。对任意 \(a\in A\) 有 \(a=p(a/p)\)，故 \(A=pA\)。
(2)\(\Leftrightarrow\)(3) 为定义等价：商群 \(A/pA\) 平凡当且仅当 \(A=pA\)。
\par
(3)\(\Rightarrow\)(1)：若 \(p\notin\mathcal P\)，则 \(A\cong\ZZ[\mathcal P^{-1}]\) 中一切元素的分母与 \(p\) 互素，从而 \(1\notin pA\)（若 \(1=pa\) 则 \(a=1/p\notin A\)），故 \(A/pA\neq 0\)。
取逆否命题即得 (3)\(\Rightarrow\)(1)。证毕。
\end{proof}

\begin{proposition}[单素数 \texorpdfstring{$p$}{p}-solenoid 上的定点计数与 Artin--Mazur \texorpdfstring{$\zeta$}{zeta} 因子]\label{prop:xi-solenoid-p-artin-mazur-zeta}
\leanverified{paper\_xi\_solenoid\_p\_artin\_mazur\_zeta}
取 \(S=\{p\}\)，令 \(A_{\{p\}}=\ZZ[1/p]\)，并记其 Pontryagin 对偶为 \(\Sigma_{\{p\}}:=\widehat{A_{\{p\}}}\)。
令 \([p]:\Sigma_{\{p\}}\to \Sigma_{\{p\}}\) 表示与对偶端乘法 \(A_{\{p\}}\xrightarrow{\times p}A_{\{p\}}\) 对应的连续群自同态。
则对每个 \(n\ge 1\)，定点集有限且满足
\[
\card{\mathrm{Fix}([p]^n)}=p^n-1.
\]
并且其 Artin--Mazur \(\zeta\) 函数
\[
\zeta_{[p]}(z)
:=\exp\!\Bigl(\sum_{n\ge 1}\frac{\card{\mathrm{Fix}([p]^n)}}{n}z^n\Bigr)
\]
在 \(|z|<1/p\) 上给出显式有理表达
\[
\zeta_{[p]}(z)=\frac{1-z}{1-pz}.
\]
\end{proposition}
\begin{proof}
定点条件 \([p]^n(x)=x\) 等价于 \((p^n-1)x=0\)，从而
\(
\mathrm{Fix}([p]^n)=\ker([p^n-1])
\)。
对偶化得
\[
\ker([m])^\vee\cong A_{\{p\}}/mA_{\{p\}}.
\]
取 \(m=p^n-1\)。
由于 \(\gcd(m,p)=1\)，在模 \(m\) 商中局部化不改变同构类型，从而自然映射诱导同构
\[
\ZZ/m\ZZ\ \xrightarrow{\ \sim\ }\ A_{\{p\}}/mA_{\{p\}}.
\]
故 \(\card{\mathrm{Fix}([p]^n)}=\card{A_{\{p\}}/mA_{\{p\}}}=m=p^n-1\)。
\par
当 \(|z|<1/p\) 时幂级数绝对收敛，并且
\[
\sum_{n\ge 1}\frac{(p^n-1)z^n}{n}
=\sum_{n\ge 1}\frac{(pz)^n}{n}-\sum_{n\ge 1}\frac{z^n}{n}
=-\log(1-pz)+\log(1-z).
\]
指数化即得 \(\zeta_{[p]}(z)=(1-z)/(1-pz)\)。证毕。
\end{proof}

\begin{corollary}[欧拉乘积的动力学实现]\label{cor:xi-artin-mazur-euler-product-zeta-ratio}
\leanverified{paper\_xi\_artin\_mazur\_euler\_product\_zeta\_ratio}
对复数 \(s\) 满足 \(\Re(s)>2\)，有绝对收敛恒等式
\[
\prod_{p}\zeta_{[p]}(p^{-s})
=\prod_{p}\frac{1-p^{-s}}{1-p^{1-s}}
=\frac{\zeta(s-1)}{\zeta(s)}.
\]
\end{corollary}
\begin{proof}
由命题 \ref{prop:xi-solenoid-p-artin-mazur-zeta}，\(\zeta_{[p]}(p^{-s})=(1-p^{-s})/(1-p^{1-s})\)。
当 \(\Re(s)>2\) 时两端欧拉乘积绝对收敛，逐项相乘合法；并且
\(
\zeta(s)=\prod_p(1-p^{-s})^{-1}
\)、\(
\zeta(s-1)=\prod_p(1-p^{1-s})^{-1}
\)。
合并即得所示恒等式。证毕。
\end{proof}

\begin{corollary}[局部因子类型的可逆性]\label{cor:xi-cdim-localization-zeta-recovers-S}
\leanverified{paper\_xi\_cdim\_localization\_zeta\_recovers\_s}
在定理 \ref{thm:xi-cdim-localization-zeta-AS} 的设定下，有限素数集 \(S\) 由 \(\zeta_{A_S}\) 唯一确定。
\end{corollary}
\begin{proof}
欧拉乘积中每个素数 \(p\) 的局部因子为
\(
(1-p^{1-s})^{-1}
\)
或
\(
(1-p^{-s})^{-1}
\)，二者作为 \(p\)-局部 Dirichlet 因子类型不同。
由素数逐点唯一分解，局部因子类型的分布唯一给出 \(S\)。
\end{proof}

\begin{theorem}[有限素数集局部化群的 \texorpdfstring{$\mathrm{Hom}$}{Hom}-范畴完全分类]\label{thm:xi-cdim-localization-hom-category-classification}
对有限素数集 \(S,T\)，定义
\[
G_S:=\ZZ[S^{-1}],
\qquad
G_T:=\ZZ[T^{-1}]
\]
并按加法群理解。则有自然同构
\[
\operatorname{Hom}(G_S,G_T)\cong
\begin{cases}
G_T,& S\subseteq T,\\[2mm]
0,& S\not\subseteq T.
\end{cases}
\]
特别地，
\[
\operatorname{End}(G_S)\cong G_S,
\qquad
\operatorname{Aut}(G_S)\cong
\Bigl\{\pm\prod_{p\in S}p^{n_p}:n_p\in\ZZ\Bigr\}.
\]
因此
\[
G_S\cong G_T
\qquad\Longleftrightarrow\qquad
S=T.
\]
\end{theorem}
\begin{proof}
任意群同态 \(\phi:G_S\to G_T\) 被 \(\phi(1)\) 唯一决定。
事实上，\(G_S\subset\QQ\) 且由元素 \(1/p^n\)（\(p\in S,\ n\ge 0\)）加法生成；由于 \(G_T\subset\QQ\) 无挠，对每个 \(p^n\) 方程
\(
p^n y=\phi(1)
\)
至多有一个解，因此 \(\phi(1/p^n)\) 一旦存在便唯一确定。

若 \(S\subseteq T\)，任取 \(q\in G_T\)，定义
\[
\phi_q(x):=qx\qquad (x\in G_S).
\]
因为 \(x\in G_S\) 的分母只含 \(S\) 中素数，而 \(S\subseteq T\)，故 \(qx\in G_T\)，从而 \(\phi_q\) 为良定义群同态。映射
\(
q\mapsto \phi_q
\)
显然为单射。反之，对任意 \(\phi\) 取 \(q:=\phi(1)\)，上段唯一性说明 \(\phi=\phi_q\)。故 \(\operatorname{Hom}(G_S,G_T)\cong G_T\)。

若 \(S\not\subseteq T\)，取 \(p\in S\setminus T\)。设 \(\phi:G_S\to G_T\) 且 \(\phi(1)=q\neq 0\)。则对每个 \(n\ge 1\)，元素 \(1/p^n\in G_S\) 的像必须满足
\[
p^n\phi(1/p^n)=q.
\]
由于 \(G_T\subset\QQ\) 且 \(p\notin T\)，元素 \(q/p^n\) 对充分大的 \(n\) 不再属于 \(G_T\)，矛盾。
因此只能有 \(q=0\)，而上段唯一性随即推出 \(\phi=0\)。故 \(\operatorname{Hom}(G_S,G_T)=0\)。

令 \(T=S\) 即得 \(\operatorname{End}(G_S)\cong G_S\)。其中 \(\phi_q\) 为自同构当且仅当 \(qG_S=G_S\)，等价于 \(q,q^{-1}\in G_S\)，也即 \(q\) 属于局部化环 \(\ZZ[S^{-1}]\) 的单位群
\[
G_S^\times=
\Bigl\{\pm\prod_{p\in S}p^{n_p}:n_p\in\ZZ\Bigr\}.
\]
最后，若 \(G_S\cong G_T\)，则存在非零同态 \(G_S\to G_T\) 与 \(G_T\to G_S\)，由上式分别推出 \(S\subseteq T\) 与 \(T\subseteq S\)，故 \(S=T\)；反向蕴含显然。证毕。
\end{proof}
\leanverified{paper\_xi\_cdim\_localization\_hom\_category\_classification}

\begin{theorem}[局部化数系包含关系的商分解与 solenoid 满射分类]\label{thm:xi-cdim-localization-quotient-solenoid-surjections}
沿用定理 \ref{thm:xi-cdim-localization-hom-category-classification} 的记号，并记
\[
\Sigma_S:=\widehat{G_S},
\qquad
\Sigma_T:=\widehat{G_T}.
\]
若 \(S\subseteq T\)，则存在规范短正合列
\[
0\longrightarrow G_S\longrightarrow G_T
\longrightarrow \bigoplus_{p\in T\setminus S} C_{p^\infty}
\longrightarrow 0
\]
以及其 Pontryagin 对偶
\[
0\longrightarrow \prod_{p\in T\setminus S}\ZZ_p
\longrightarrow \Sigma_T
\xrightarrow{\ \pi_{T\to S}\ }\Sigma_S
\longrightarrow 0.
\]
此外，连续满同态 \(\Sigma_T\twoheadrightarrow \Sigma_S\) 存在当且仅当 \(S\subseteq T\)。

更强地，当 \(S\subseteq T\) 时，任意连续满同态 \(f:\Sigma_T\twoheadrightarrow \Sigma_S\) 都可写成
\[
f=[m]\circ \pi_{T\to S}\circ \alpha
\]
的形式，其中 \([m]:\Sigma_S\to\Sigma_S\) 表示乘以 \(m\) 的连续自同态，\(\alpha\in\operatorname{Aut}(\Sigma_T)\)，且 \(m\ge 1\) 的所有素因子均不属于 \(T\)。
等价地，其核满足
\[
\ker(f)\cong
\left(\prod_{p\in T\setminus S}\ZZ_p\right)\times \ZZ/m\ZZ.
\]
特别地，规范满同态 \(\pi_{T\to S}\) 对应于 \(m=1\)，其核恰为
\(
\prod_{p\in T\setminus S}\ZZ_p
\)。
\end{theorem}
\leanverified{paper\_xi\_cdim\_localization\_quotient\_solenoid\_surjections}
\begin{proof}
先证明规范商分解。对任意有限素数集 \(U\)，群 \(G_U/\ZZ\) 可识别为 \(\QQ/\ZZ\) 中“分母只允许落在 \(U\) 内”的子群，因此
\[
G_U/\ZZ\cong \bigoplus_{p\in U} C_{p^\infty}.
\]
当 \(S\subseteq T\) 时，
\[
G_T/G_S\cong (G_T/\ZZ)/(G_S/\ZZ)\cong \bigoplus_{p\in T\setminus S} C_{p^\infty},
\]
即得第一条短正合列。对偶化并使用
\(
\widehat{C_{p^\infty}}\cong \ZZ_p
\)
便得到第二条短正合列。

若 \(S\subseteq T\)，规范投影 \(\pi_{T\to S}\) 已给出一个连续满同态。
反之，若存在连续满同态 \(f:\Sigma_T\twoheadrightarrow \Sigma_S\)，则对偶化得到单射
\[
f^\vee:G_S\hookrightarrow G_T.
\]
由定理 \ref{thm:xi-cdim-localization-hom-category-classification}，这只有在 \(S\subseteq T\) 时才可能。故满同态的存在性当且仅当 \(S\subseteq T\)。

现在设 \(S\subseteq T\)，并取任意满同态 \(f:\Sigma_T\twoheadrightarrow \Sigma_S\)。
其对偶单射 \(f^\vee:G_S\hookrightarrow G_T\) 由同一定理必为某个非零 \(q\in G_T\) 所诱导的乘法映射 \(\phi_q(x)=qx\)。
把 \(q\) 写成
\[
q=u\,m,
\]
其中 \(u\in G_T^\times\) 为单位，且 \(m\ge 1\) 的所有素因子均不属于 \(T\)。
于是
\[
\phi_q=(\times u)\circ \iota\circ (\times m),
\]
其中 \(\iota:G_S\hookrightarrow G_T\) 为规范包含。
对偶化即得
\[
f=[m]\circ \pi_{T\to S}\circ \alpha,
\]
其中 \(\alpha:=\widehat{(\times u)}\in\operatorname{Aut}(\Sigma_T)\)。

最后计算核。由于 \(mG_S\subseteq G_S\subseteq G_T\)，有短正合列
\[
0\longrightarrow G_S/mG_S\longrightarrow G_T/mG_S\longrightarrow G_T/G_S\longrightarrow 0.
\]
因 \(m\) 的所有素因子都不属于 \(T\)（从而也不属于 \(S\)），局部化不改变模 \(m\) 商，故
\[
G_S/mG_S\cong \ZZ/m\ZZ.
\]
另一方面，\(G_T/G_S\cong \bigoplus_{p\in T\setminus S} C_{p^\infty}\) 为可除挠群，因而作为阿贝尔群是内射对象；上式分裂，从而
\[
\operatorname{coker}(f^\vee)\cong G_T/qG_S\cong \ZZ/m\ZZ\oplus \bigoplus_{p\in T\setminus S} C_{p^\infty}.
\]
再取 Pontryagin 对偶便得到
\[
\ker(f)\cong
\left(\prod_{p\in T\setminus S}\ZZ_p\right)\times \ZZ/m\ZZ.
\]
证毕。
\end{proof}

\begin{theorem}[有限素数局部化 solenoid 之间的连续映射完全分类]\label{thm:xi-cdim-localization-solenoid-continuous-hom-classification}
\leanverified{paper\_xi\_cdim\_localization\_solenoid\_continuous\_hom\_classification}
沿用定理 \ref{thm:xi-cdim-localization-hom-category-classification} 与
\ref{thm:xi-cdim-localization-quotient-solenoid-surjections} 的记号，记
\[
\Sigma_S:=\widehat{G_S},
\qquad
\Sigma_T:=\widehat{G_T}.
\]
则存在自然同构
\[
\operatorname{Hom}_{\mathrm{cts}}(\Sigma_S,\Sigma_T)
\cong
\operatorname{Hom}(G_T,G_S)
\cong
\begin{cases}
G_S,& T\subseteq S,\\[1mm]
0,& T\not\subseteq S.
\end{cases}
\]
其中右端同构按加法群理解。

特别地，当 \(T\subseteq S\) 时，每个非零连续同态
\(
\Sigma_S\to \Sigma_T
\)
都是满射。
\end{theorem}
\begin{proof}
Pontryagin 对偶把紧交换群范畴与离散交换群范畴反变等价，因此有自然同构
\[
\operatorname{Hom}_{\mathrm{cts}}(\Sigma_S,\Sigma_T)
\cong
\operatorname{Hom}(G_T,G_S).
\]
将定理 \ref{thm:xi-cdim-localization-hom-category-classification} 应用于有序对 \((T,S)\)，便得
\[
\operatorname{Hom}(G_T,G_S)\cong
\begin{cases}
G_S,& T\subseteq S,\\
0,& T\not\subseteq S.
\end{cases}
\]
这就给出所述分类。

若 \(T\subseteq S\) 且 \(0\neq q\in G_S\)，则对应离散端同态
\[
\phi_q:G_T\to G_S,\qquad x\mapsto qx
\]
为单射；因为 \(G_S\subset\QQ\) 无挠，而 \(q\neq 0\) 时 \(qx=0\) 只能推出 \(x=0\)。
对偶化后即得对应的连续同态 \(\widehat{\phi_q}:\Sigma_S\to\Sigma_T\) 为满射。
\end{proof}

\begin{corollary}[有限素数局部化 solenoid 的满商反变序与连续端同态环]\label{cor:xi-cdim-localization-solenoid-quotient-order-endomorphism}
\leanverified{paper\_xi\_cdim\_localization\_solenoid\_quotient\_order\_endomorphism}
沿用定理 \ref{thm:xi-cdim-localization-solenoid-continuous-hom-classification} 的记号。
则成立：
\begin{enumerate}
  \item 存在连续满同态
  \[
  \Sigma_S\twoheadrightarrow \Sigma_T
  \]
  当且仅当 \(T\subseteq S\)。
  因而在同构类层面，有限素数子集的包含格
  \(
  (\mathcal P_f(\mathbb P),\subseteq)
  \)
  与局部化 solenoid 族的满商偏序之间形成严格反变对应。
  \item 连续端同态环满足
  \[
  \operatorname{End}_{\mathrm{cts}}(\Sigma_S)\cong \ZZ[S^{-1}]
  \]
  （作为交换环），而连续自同构群满足
  \[
  \operatorname{Aut}_{\mathrm{cts}}(\Sigma_S)
  \cong
  \ZZ[S^{-1}]^\times
  =
  \left\{\pm\prod_{p\in S}p^{n_p}: n_p\in\ZZ\right\}.
  \]
\end{enumerate}
\end{corollary}
\begin{proof}
第一条由定理 \ref{thm:xi-cdim-localization-solenoid-continuous-hom-classification} 直接得到：若 \(T\subseteq S\)，任取 \(0\neq q\in G_S\) 即给出一条满射 \(\Sigma_S\twoheadrightarrow \Sigma_T\)；反之若存在连续满射，则其对偶为单射
\(
G_T\hookrightarrow G_S
\)，于是必有 \(T\subseteq S\)。

对第二条，Pontryagin 对偶把端同态环识别为离散端端同态环的反环：
\[
\operatorname{End}_{\mathrm{cts}}(\Sigma_S)\cong \operatorname{End}(G_S)^{\mathrm{op}}.
\]
由定理 \ref{thm:xi-cdim-localization-hom-category-classification}，
\(
\operatorname{End}(G_S)\cong G_S
\)
并由乘法映射 \(m_q(x)=qx\) 实现。
又
\[
m_q\circ m_r=m_{qr},
\]
故 \(\operatorname{End}(G_S)\) 与局部化环 \(\ZZ[S^{-1}]\) 同构。
由于该环交换，其反环与自身一致，遂得
\[
\operatorname{End}_{\mathrm{cts}}(\Sigma_S)\cong \ZZ[S^{-1}].
\]
可逆元恰为局部化环的单位群，因此
\[
\operatorname{Aut}_{\mathrm{cts}}(\Sigma_S)\cong \ZZ[S^{-1}]^\times
=
\left\{\pm\prod_{p\in S}p^{n_p}: n_p\in\ZZ\right\}.
\]
证毕。
\end{proof}

\begin{theorem}[有限素数局部化 solenoid 的圆周商扩张在非平凡情形下不分裂]\label{thm:xi-cdim-localization-solenoid-circle-extension-nonsplit}
\leanverified{paper\_xi\_cdim\_localization\_solenoid\_circle\_extension\_nonsplit}
设 \(S\neq\varnothing\) 为有限素数集。
则定理 \ref{thm:cdim-star-z1s-dual-extension} 中的短正合列
\[
0\longrightarrow \prod_{p\in S}\ZZ_p
\longrightarrow \Sigma_S
\longrightarrow \TT
\longrightarrow 0
\]
在拓扑群范畴中不分裂。
等价地，
\[
\Sigma_S\not\cong \TT\times \prod_{p\in S}\ZZ_p
\qquad (S\neq\varnothing).
\]
\end{theorem}
\begin{proof}
若该短正合列分裂，则存在连续截面 \(\sigma:\TT\to\Sigma_S\)，从而
\[
\Sigma_S\cong \TT\times \prod_{p\in S}\ZZ_p.
\]
对偶化得到离散交换群的分裂短正合列
\[
0\longrightarrow \ZZ
\longrightarrow G_S
\longrightarrow \bigoplus_{p\in S}C_{p^\infty}
\longrightarrow 0,
\]
从而
\[
G_S\cong \ZZ\oplus \bigoplus_{p\in S}C_{p^\infty}.
\]
但当 \(S\neq\varnothing\) 时，右端含有非零挠元；另一方面
\[
G_S=\ZZ[S^{-1}]\subset \QQ
\]
作为 \(\QQ\) 的加法子群是无挠群，矛盾。
故该扩张不能分裂。
\end{proof}

\paragraph{有限维环面观测与 prime-register 可见性边界}
在有限素数局部化群与其对偶 solenoid 的 \(\mathrm{Hom}\)-范畴、满商偏序与不分裂扩张已经完全刻画之后，还可进一步把“有限维环面可见量”的可实现边界压缩为下列四条结论。

\begin{theorem}[有限维环面既不能容纳也不能生成局部化 solenoid]\label{thm:xi-cdim-localization-solenoid-not-subgroup-or-quotient-of-finite-torus}
\leanverified{paper\_xi\_cdim\_localization\_solenoid\_not\_subgroup\_or\_quotient\_of\_finite\_torus}
设 \(S\subset\mathbb P\) 为非空有限素数集，并记
\[
G_S:=\ZZ[S^{-1}],
\qquad
\Sigma_S:=\widehat{G_S}.
\]
则对任意整数 \(n\ge 1\)：
\begin{enumerate}
  \item 不存在连续单射
  \[
  \Sigma_S\hookrightarrow \TT^n.
  \]
  \item 不存在连续满射
  \[
  \TT^n\twoheadrightarrow \Sigma_S.
  \]
\end{enumerate}
\end{theorem}
\begin{proof}
先说明 \(G_S\) 在 \(S\neq\varnothing\) 时不是有限生成阿贝尔群。任取 \(p\in S\)。若 \(H\subset G_S\) 由有限集合
\[
\left\{\frac{a_1}{N},\dots,\frac{a_r}{N}\right\}
\]
生成，则必有 \(H\subseteq \frac1N\ZZ\)。取 \(k\) 足够大使 \(p^k\nmid N\)，则
\[
\frac{1}{p^k}\notin \frac1N\ZZ,
\]
故 \(H\) 不可能包含全部 \(p^{-k}\)。因此 \(G_S\) 不是有限生成群。

现证第 \(1\) 条。若存在连续单射
\[
\iota:\Sigma_S\hookrightarrow \TT^n,
\]
则因 \(\Sigma_S\) 紧、\(\TT^n\) Hausdorff，\(\iota\) 为到其闭像的拓扑同构。对偶化得到满射
\[
\widehat\iota:\widehat{\TT^n}\cong \ZZ^n\twoheadrightarrow \widehat{\Sigma_S}\cong G_S.
\]
但 \(\ZZ^n\) 的任何商群都是有限生成阿贝尔群，这与 \(G_S\) 非有限生成矛盾。

再证第 \(2\) 条。若存在连续满射
\[
\pi:\TT^n\twoheadrightarrow \Sigma_S,
\]
则 Pontryagin 对偶给出单射
\[
\widehat\pi:G_S\hookrightarrow \widehat{\TT^n}\cong \ZZ^n.
\]
而 \(\ZZ^n\) 的任意子群仍然有限生成，与 \(G_S\) 非有限生成矛盾。证毕。
\end{proof}

\begin{theorem}[所有有限维环面值观测都因子化为单一相位通道]\label{thm:xi-cdim-localization-solenoid-torus-observation-single-phase-factorization}
\leanverified{paper\_xi\_cdim\_localization\_solenoid\_torus\_observation\_single\_phase\_factorization}
沿用定理 \ref{thm:xi-cdim-localization-solenoid-not-subgroup-or-quotient-of-finite-torus} 的记号。对任意连续群同态
\[
f:\Sigma_S\longrightarrow \TT^n,
\]
存在唯一的极小素数子集 \(T_f\subseteq S\) 与一个连续同态
\[
\widetilde f:\Sigma_{T_f}\longrightarrow \TT^n
\]
使得
\[
f=\widetilde f\circ \pi_{S\to T_f},
\]
其中 \(\pi_{S\to T_f}:\Sigma_S\twoheadrightarrow \Sigma_{T_f}\) 为自然满射。并且 \(\widetilde f\) 进一步因子化经过一个圆：
\[
\Sigma_{T_f}\twoheadrightarrow \TT\xrightarrow{\ \alpha\ }\TT^n.
\]
换言之，任意有限维环面值观测至多只看到一条 rank-\(1\) 相位方向。
\end{theorem}
\begin{proof}
对偶化 \(f\) 得到群同态
\[
\widehat f:\ZZ^n\longrightarrow G_S.
\]
其像
\[
H:=\operatorname{im}(\widehat f)
\]
是 \(G_S\) 的有限生成子群。任取 \(H\) 的一组生成元
\[
\frac{a_1}{b_1},\dots,\frac{a_r}{b_r}\in G_S\subset\QQ,
\]
令 \(N:=\operatorname{lcm}(b_1,\dots,b_r)\)。则
\[
H\subseteq \frac1N\ZZ.
\]
由于 \(\frac1N\ZZ\) 作为加法群同构于 \(\ZZ\)，故 \(H\) 必为循环群。若 \(H=0\)，则取
\[
T_f=\varnothing,\qquad \widetilde f=0:\Sigma_{\varnothing}=\TT\to \TT^n,
\]
便已得到所需分解。下设 \(H\neq 0\)。

取 \(a\in G_S\) 使
\[
H=\ZZ a.
\]
把 \(a\) 写成既约分数 \(a=u/v\)（\(u\in\ZZ\)、\(v\in\NN\)、\(\gcd(u,v)=1\)），并定义 \(T_f\) 为 \(v\) 的素数支撑。则 \(a\in G_{T_f}\)，并且若 \(T\subseteq S\) 满足 \(H\subseteq G_T\)，则必有 \(a\in G_T\)，从而 \(T_f\subseteq T\)。故 \(T_f\) 是唯一极小的可行素数子集。

由 \(H=\ZZ a\) 可取群同态
\[
\psi:\ZZ^n\longrightarrow \ZZ
\]
与
\[
\mu_a:\ZZ\longrightarrow G_{T_f},
\qquad
k\longmapsto ka
\]
使得
\[
\widehat f
=
\iota_{T_f,S}\circ \mu_a\circ \psi,
\]
其中 \(\iota_{T_f,S}:G_{T_f}\hookrightarrow G_S\) 为自然包含。
对偶化得到
\[
f
=
\widehat\psi\circ \widehat{\mu_a}\circ \widehat{\iota_{T_f,S}}
=
\alpha\circ \chi_a\circ \pi_{S\to T_f},
\]
其中
\[
\pi_{S\to T_f}:=\widehat{\iota_{T_f,S}}:\Sigma_S\twoheadrightarrow \Sigma_{T_f},
\qquad
\chi_a:=\widehat{\mu_a}:\Sigma_{T_f}\to \TT,
\qquad
\alpha:=\widehat\psi:\TT\to \TT^n.
\]
由于 \(a\neq 0\) 且 \(G_{T_f}\subset\QQ\) 无挠，\(\mu_a\) 为单射；因此其对偶 \(\chi_a\) 为满射。
令
\[
\widetilde f:=\alpha\circ \chi_a:\Sigma_{T_f}\to \TT^n,
\]
便得到所述分解。证毕。
\end{proof}

\begin{corollary}[局部化 solenoid 的一切有限维酉表示都来自圆]\label{cor:xi-cdim-localization-solenoid-finite-dimensional-unitary-representations-from-circle}
设
\[
\rho:\Sigma_S\longrightarrow \mathrm{U}(N)
\]
为连续有限维酉表示。则：
\begin{enumerate}
  \item \(\rho\) 完全可约为一维角色之直和；
  \item 每个角色都因子化经过某个圆商；
  \item 整个表示 \(\rho\) 也因子化经过一个圆：
  \[
  \Sigma_S\twoheadrightarrow \TT\xrightarrow{\ \beta\ }\mathrm{U}(N).
  \]
\end{enumerate}
特别地，当 \(S\neq\varnothing\) 时，\(\Sigma_S\) 不存在忠实有限维酉表示。
\end{corollary}
\begin{proof}
由于 \(\Sigma_S\) 为紧交换群，其连续有限维酉表示可同时酉对角化，故
\[
\rho\cong \chi_1\oplus\cdots\oplus\chi_N
\]
为一维角色之直和。令
\[
f:=(\chi_1,\dots,\chi_N):\Sigma_S\to \TT^N.
\]
由定理 \ref{thm:xi-cdim-localization-solenoid-torus-observation-single-phase-factorization}，\(f\) 因子化为
\[
\Sigma_S\twoheadrightarrow \TT\xrightarrow{\ \alpha\ }\TT^N.
\]
把 \(\TT^N\) 的对角表示复合到 \(\mathrm{U}(N)\) 中，便得到
\[
\rho=\beta\circ q
\]
的分解，其中 \(q:\Sigma_S\twoheadrightarrow \TT\) 为连续满射，\(\beta:\TT\to \mathrm{U}(N)\) 为连续酉表示。故前三条成立。

若 \(\rho\) 忠实，则其像是 \(\mathrm{U}(N)\) 的紧交换子群，因而经酉共轭后包含于标准对角环面 \(\TT^N\)。这便给出连续单射
\[
\Sigma_S\hookrightarrow \TT^N,
\]
与定理 \ref{thm:xi-cdim-localization-solenoid-not-subgroup-or-quotient-of-finite-torus} 矛盾。故不存在忠实有限维酉表示。证毕。
\end{proof}

\begin{theorem}[受限分母支撑的角色族恰好看见对应的 prime-register 商]\label{thm:xi-cdim-localization-solenoid-character-family-sees-exact-prime-register-quotient}
对任意子集 \(T\subseteq S\)，定义角色族
\[
\mathcal X_T
:=
\{\chi_a:\Sigma_S\to \TT:\ a\in G_T\subseteq G_S\},
\]
其中 \(a\in G_S\) 与角色 \(\chi_a\) 通过 Pontryagin 对偶一一对应。则有精确恒等式
\[
\boxed{\;
\bigcap_{\chi\in\mathcal X_T}\ker(\chi)
=
\ker(\pi_{S\to T})
\cong
\prod_{p\in S\setminus T}\ZZ_p.\;}
\]
换言之，只允许分母素支撑落在 \(T\) 内的一维相位观测时，它们恰好忽略补集 \(S\setminus T\) 的全部 \(p\)-进寄存器，而不会忽略更多，也不会忽略更少。
\end{theorem}
\begin{proof}
由包含 \(G_T\hookrightarrow G_S\) 的对偶化，得到自然满射
\[
\pi_{S\to T}:\Sigma_S\twoheadrightarrow \Sigma_T.
\]
将定理 \ref{thm:xi-cdim-localization-quotient-solenoid-surjections} 应用于有序对 \((T,S)\)，可得短正合列
\[
0\longrightarrow \prod_{p\in S\setminus T}\ZZ_p
\longrightarrow \Sigma_S
\xrightarrow{\ \pi_{S\to T}\ }
\Sigma_T
\longrightarrow 0.
\]
故
\[
\ker(\pi_{S\to T})\cong \prod_{p\in S\setminus T}\ZZ_p.
\]

另一方面，\(\mathcal X_T\) 正是 \(\Sigma_T\) 上全体角色经 \(\pi_{S\to T}\) 拉回所得的角色族。由于紧交换群的全体角色分离点，
\[
\bigcap_{\chi\in \widehat{\Sigma_T}}\ker(\chi)=\{1\}.
\]
把这一等式经由 \(\pi_{S\to T}\) 拉回，便得到
\[
\bigcap_{\chi\in\mathcal X_T}\ker(\chi)=\ker(\pi_{S\to T}).
\]
与上一步合并即得所示公式。证毕。
\end{proof}

\input{para__xi-cdim-localization-solenoid-unitary-normalform-torus-input-certificate-gap}

\endinput


\endinput


\begin{theorem}[由素数账本 \texorpdfstring{$|G/p^kG|$}{|G/p^kG|} 完整重构有限生成交换群]\label{thm:xi-cdim-prime-ledger-complete-reconstruction}
设 \(G,H\) 为有限生成离散交换群。若对所有素数 \(p\) 与所有 \(k\ge 1\) 都有
\[
|G/p^kG|=|H/p^kH|,
\]
则 \(G\cong H\)。
特别地，若对所有 \(n\ge 1\) 有 \(|G/nG|=|H/nH|\)，亦推出 \(G\cong H\)。
\end{theorem}
\begin{proof}
写
\[
G\cong \ZZ^{r_G}\oplus F,\qquad H\cong \ZZ^{r_H}\oplus F',
\]
其中 \(F,F'\) 有限。
定义
\[
e_G(p):=\log_p|G/pG|,\qquad e_H(p):=\log_p|H/pH|.
\]
由
\[
|G/pG|=p^{r_G}|F/pF|,\qquad |H/pH|=p^{r_H}|F'/pF'|
\]
知 \(e_G(p)\ge r_G\)、\(e_H(p)\ge r_H\)，且当 \(p\nmid |F|\)（分别 \(p\nmid |F'|\)）时取等。
故
\[
r_G=\min_p e_G(p),\qquad r_H=\min_p e_H(p).
\]
由假设 \(e_G(p)=e_H(p)\) 对所有 \(p\) 成立，得 \(r_G=r_H=:r\)。

固定素数 \(p\)，定义
\[
b_{G,p}(k):=v_p(|G/p^kG|)-kr,\qquad
b_{H,p}(k):=v_p(|H/p^kH|)-kr.
\]
把 \(F\) 的 \(p\)-primary 分量写为
\[
F_p\cong \bigoplus_{i=1}^{t_p}\ZZ/p^{\alpha_i}\ZZ,\qquad
\alpha_1\ge\cdots\ge\alpha_{t_p}\ge 1.
\]
则
\[
b_{G,p}(k)=v_p(|F_p/p^kF_p|)=\sum_{i=1}^{t_p}\min(\alpha_i,k).
\]
令
\[
m_{G,p}(k):=b_{G,p}(k)-b_{G,p}(k-1)=\#\{i:\alpha_i\ge k\},
\]
则
\[
\#\{i:\alpha_i=k\}=m_{G,p}(k)-m_{G,p}(k+1).
\]
因此序列 \(\{b_{G,p}(k)\}_{k\ge 1}\) 唯一决定 \(F_p\) 的不变量因子；同理对 \(H\) 成立。
由假设 \(b_{G,p}(k)=b_{H,p}(k)\)，得所有 \(p\)-primary 分量一致，故 \(F\cong F'\)。
再结合 \(r_G=r_H\)，即得 \(G\cong H\)。

若 \(|G/nG|=|H/nH|\) 对所有 \(n\) 成立，则取 \(n=p^k\) 即回到上式条件。证毕。
\end{proof}
\leanverified{paper\_xi\_cdim\_prime\_ledger\_complete\_reconstruction}

\begin{proposition}[好素数一次观测给出圆维，坏素数支撑给出有限账本]\label{prop:xi-cdim-good-prime-readout-and-bad-support}
设 \(G\cong \ZZ^r\oplus F\)（\(r=\cdim(G)\)），并记
\[
S(G):=\{p:\ p\mid |F|\}.
\]
则：
\begin{enumerate}
  \item 对任意 \(p\notin S(G)\) 与任意 \(k\ge 1\)，
  \[
  |G/p^kG|=p^{kr}.
  \]
  \item 若某素数 \(p\) 满足 \(|G/pG|>p^r\)，则 \(p\in S(G)\)。
\end{enumerate}
\end{proposition}
\begin{proof}
若 \(p\nmid |F|\)，则乘以 \(p^k\) 在 \(F\) 上可逆，故 \(F/p^kF\) 平凡，于是
\[
|G/p^kG|=|(\ZZ/p^k\ZZ)^r|\,|F/p^kF|=p^{kr}.
\]
反过来，若 \(p\mid |F|\)，则 \(F/pF\) 非平凡，故
\[
|G/pG|=p^r|F/pF|>p^r.
\]
命题两条即得。证毕。
\end{proof}
\leanverified{paper\_xi\_cdim\_good\_prime\_readout\_and\_bad\_support}

\paragraph{Gödel--prime 外幂盒体的 Löwner 椭球与体--界全息权重}\label{par:xi-godel-loewner-exterior-volume-holography}
本段给出一条与 prime-register/Gödel 外置口径兼容的几何接口：把外幂层的乘法边界数据组织为一族坐标盒体，并以其最小体积包含椭球（Löwner 椭球）为“体量”读出；其对数体积对边界总和呈现严格线性权重，且可由两阶体量作显式反演。

\begin{theorem}[外幂盒体的 Löwner 椭球与严格线性全息权重]\label{thm:xi-godel-loewner-exterior-box-loewner-holographic-weight}
设 \(k\ge 1\)、\(1\le t\le k\)，并取向量 \(L=(L_1,\dots,L_k)\) 满足 \(L_i>0\)。
记
\[
\mathcal I_{k,t}:=\{I\subseteq\{1,\dots,k\}:\ |I|=t\},
\qquad
N:=|\mathcal I_{k,t}|=\binom{k}{t},
\]
并对每个 \(I\in\mathcal I_{k,t}\) 定义 \(L_I:=\prod_{i\in I}L_i\)。
考虑坐标盒体
\[
K_{k,t}(L):=\prod_{I\in\mathcal I_{k,t}}[0,L_I]\ \subset\ \RR^{N}.
\]
令 \(E_{k,t}(L)\) 为包含 \(K_{k,t}(L)\) 的最小体积椭球（Löwner 椭球），并记 \(\kappa_N\) 为 \(\RR^N\) 中单位球体积。
则成立：
\begin{enumerate}
  \item \(E_{k,t}(L)\) 为坐标对齐椭球，其中心为 \(\bigl(\tfrac{L_I}{2}\bigr)_{I\in\mathcal I_{k,t}}\)，半轴长为 \(\sqrt N\,\tfrac{L_I}{2}\)（\(I\in\mathcal I_{k,t}\)）。
  \item 其体积满足闭式
  \[
  \operatorname{vol}\bigl(E_{k,t}(L)\bigr)
  =
  \kappa_N\,N^{N/2}\,2^{-N}\,\prod_{I\in\mathcal I_{k,t}}L_I
  =
  \kappa_N\,N^{N/2}\,2^{-N}\,
  \Bigl(\prod_{i=1}^{k}L_i\Bigr)^{\binom{k-1}{t-1}}.
  \]
  \item 单位维度对数体积给出严格线性权重恒等式
  \[
  \frac1N\log\operatorname{vol}\bigl(E_{k,t}(L)\bigr)
  =
  \frac1N\log\kappa_N+\frac12\log N-\log 2+\frac{t}{k}\sum_{i=1}^{k}\log L_i.
  \]
\end{enumerate}
\end{theorem}
\leanverified{paper\_xi\_godel\_loewner\_exterior\_box\_loewner\_holographic\_weight}
\begin{proof}
将 \(K_{k,t}(L)\) 平移为中心对称体
\[
\widetilde K_{k,t}(L):=\prod_{I\in\mathcal I_{k,t}}\Bigl[-\frac{L_I}{2},\frac{L_I}{2}\Bigr].
\]
平移不改变最小体积包含椭球的体积与形状，仅改变中心，故只需处理 \(\widetilde K_{k,t}(L)\)。
\par
先对立方体 \([-1,1]^N\) 说明：其对称群包含超正交群（符号翻转与坐标置换）作用。
对任意包含 \([-1,1]^N\) 的椭球 \(E\)，对该有限群平均其二次型（保持体积不增）得到一个同样包含 \([-1,1]^N\) 且在该群下不变的椭球 \(E^{\mathrm{sym}}\)。
群不变性强制 \(E^{\mathrm{sym}}\) 为欧氏球 \(B_R\)。
又因顶点 \(v=(1,\dots,1)\) 满足 \(\|v\|_2=\sqrt N\)，任意包含 \([-1,1]^N\) 的球半径必满足 \(R\ge \sqrt N\)，故最小体积包含椭球为 \(B_{\sqrt N}\)。
\par
对一般盒体 \(\widetilde K_{k,t}(L)\)，取对角线性变换
\[
A:=\mathrm{diag}\Bigl(\frac{L_I}{2}\Bigr)_{I\in\mathcal I_{k,t}},
\]
则 \(A([-1,1]^N)=\widetilde K_{k,t}(L)\)。
任意包含 \(\widetilde K_{k,t}(L)\) 的椭球经 \(A^{-1}\) 拉回即为包含 \([-1,1]^N\) 的椭球，体积按 \(|\det A|\) 缩放，故最小体积包含椭球为 \(A(B_{\sqrt N})\)；
其半轴长为 \(\sqrt N\,\tfrac{L_I}{2}\)，并得到体积闭式
\[
\operatorname{vol}\bigl(E_{k,t}(L)\bigr)=\kappa_N\prod_{I\in\mathcal I_{k,t}}\Bigl(\sqrt N\,\frac{L_I}{2}\Bigr)
=\kappa_N\,N^{N/2}\,2^{-N}\,\prod_I L_I.
\]
\par
最后，乘积恒等式
\[
\prod_{I\in\mathcal I_{k,t}}L_I
=
\prod_{I}\prod_{i\in I}L_i
=
\prod_{i=1}^{k}L_i^{|\{I\in\mathcal I_{k,t}:\ i\in I\}|}
=
\prod_{i=1}^{k}L_i^{\binom{k-1}{t-1}}
\]
由计数 \(|\{I:\ i\in I\}|=\binom{k-1}{t-1}\) 得出。
对体积公式取对数并除以 \(N\)，再用
\(
\binom{k-1}{t-1}/\binom{k}{t}=t/k
\)
即得线性权重恒等式。证毕。
\end{proof}

\begin{theorem}\label{thm:xi-local-exterior-box-loewner-tomography-inversion}
设 \(k\ge 2\) 且 \(1\le t\le k-1\)，并取向量 \(L=(L_1,\dots,L_k)\) 满足 \(L_i>0\)。
对每个 \(i\in\{1,\dots,k\}\) 定义局部指标族与维数
\[
\mathcal I_{k,t}(i):=\{I\subseteq\{1,\dots,k\}:\ |I|=t,\ i\in I\},
\qquad
N:=|\mathcal I_{k,t}(i)|=\binom{k-1}{t-1}.
\]
并对每个 \(I\subseteq\{1,\dots,k\}\) 记 \(L_I:=\prod_{j\in I}L_j\)。
考虑局部外幂盒体
\[
K_{k,t}^{(i)}(L):=\prod_{I\in \mathcal I_{k,t}(i)}[0,L_I]\ \subset\ \RR^{N},
\]
并令 \(E_{k,t}^{(i)}(L)\) 表示包含 \(K_{k,t}^{(i)}(L)\) 的最小体积外接椭球（Löwner 椭球），\(\kappa_N:=\operatorname{vol}(B_2^N)\)。
则向量
\(
\bigl(\operatorname{vol}(E_{k,t}^{(1)}(L)),\dots,\operatorname{vol}(E_{k,t}^{(k)}(L))\bigr)
\)
唯一确定 \(L\)，并且存在显式反演式：令
\[
R_i:=\log\operatorname{vol}(E_{k,t}^{(i)}(L))-\log\kappa_N-\frac{N}{2}\log N+N\log 2,
\qquad (1\le i\le k),
\]
并记
\[
\alpha:=\binom{k-2}{t-1},
\qquad
\beta:=\binom{k-2}{t-2}.
\]
则对每个 \(i\in\{1,\dots,k\}\) 有
\[
\log L_i
=\frac{1}{\alpha}R_i-\frac{\beta}{\alpha(\alpha+k\beta)}\sum_{j=1}^k R_j,
\]
从而 \(L_i=\exp(\log L_i)\) 可由局部 Löwner 体量一次性复原。
\leanverified{paper\_xi\_local\_exterior\_box\_loewner\_tomography\_inversion}
\end{theorem}
\begin{proof}
对轴对齐盒体 \(\prod_{m=1}^{N}[0,a_m]\subset\RR^N\)，其 Löwner 椭球为对角线性像 \(A(B_2^N(\sqrt N))\)（\(A=\mathrm{diag}(a_m/2)\)），故
\[
\log\operatorname{vol}(E)
=\log\kappa_N+\frac N2\log N-N\log2+\sum_{m=1}^{N}\log a_m.
\]
应用于 \(K_{k,t}^{(i)}(L)\) 得
\[
R_i=\sum_{I\in\mathcal I_{k,t}(i)}\log L_I.
\]
而
\[
\sum_{I\in\mathcal I_{k,t}(i)}\log L_I
=\sum_{I\ni i}\ \sum_{j\in I}\log L_j
=\binom{k-1}{t-1}\log L_i+\binom{k-2}{t-2}\sum_{j\ne i}\log L_j
=\alpha\log L_i+\beta\sum_{j=1}^k\log L_j,
\]
其中最后一步使用
\(
\binom{k-1}{t-1}-\binom{k-2}{t-2}=\binom{k-2}{t-1}=\alpha
\)。
于是 \(R=(R_i)_{i=1}^k\) 与 \(x=(\log L_i)_{i=1}^k\) 满足线性关系
\[
R=(\alpha I_k+\beta J_k)x,
\qquad
J_k:=\mathbf 1\mathbf 1^{\mathsf T}.
\]
因 \(\alpha>0\) 且 \(\alpha+k\beta>0\)，矩阵 \(\alpha I_k+\beta J_k\) 可逆，并且
\[
(\alpha I_k+\beta J_k)^{-1}
=\frac{1}{\alpha}I_k-\frac{\beta}{\alpha(\alpha+k\beta)}J_k.
\]
右乘 \(R\) 即得所述反演式，从而 \(L\) 由 \(\{\operatorname{vol}(E_{k,t}^{(i)}(L))\}_{i=1}^k\) 唯一确定。证毕。
\end{proof}

\begin{corollary}[两阶体量的全息反演：边界对数和由体积唯一确定]\label{cor:xi-godel-loewner-two-level-holographic-inversion}
在定理 \ref{thm:xi-godel-loewner-exterior-box-loewner-holographic-weight} 的记号下，取任意 \(1\le t_1<t_2\le k\)，令 \(N_j:=\binom{k}{t_j}\)。
则边界量
\(
S(L):=\sum_{i=1}^{k}\log L_i
\)
可由两阶 Löwner 体量唯一确定，并满足显式反演式
\[
S(L)=\frac{k}{t_1-t_2}\left[
\frac{1}{N_1}\log\operatorname{vol}\bigl(E_{k,t_1}(L)\bigr)-\frac{1}{N_2}\log\operatorname{vol}\bigl(E_{k,t_2}(L)\bigr)
-\Big(\frac{1}{N_1}\log\kappa_{N_1}-\frac{1}{N_2}\log\kappa_{N_2}\Big)
-\frac12\log\frac{N_1}{N_2}
\right].
\]
\end{corollary}
\begin{proof}
对 \(t=t_1,t_2\) 将定理 \ref{thm:xi-godel-loewner-exterior-box-loewner-holographic-weight}(3) 两式相减，未知量仅为 \(S(L)\)，解得所示反演式。证毕。
\end{proof}

\begin{corollary}[Gödel 整数的体量译码：两阶 Löwner 体量确定乘法外置]\label{cor:xi-godel-integer-decoding-from-two-loewner-volumes}
\leanverified{paper\_xi\_godel\_integer\_decoding\_from\_two\_loewner\_volumes}
取互异素数 \(p_1,\dots,p_k\) 与指数向量 \(a=(a_1,\dots,a_k)\in\ZZ_{\ge 0}^k\)，并令
\[
L_i:=p_i^{a_i}\quad(1\le i\le k),
\qquad
N:=\prod_{i=1}^{k}p_i^{a_i}\in\ZZ_{\ge 1}.
\]
在推论 \ref{cor:xi-godel-loewner-two-level-holographic-inversion} 的设定下，任取 \(1\le t_1<t_2\le k\)，两阶体量
\(
\operatorname{vol}(E_{k,t_1}(L))
\)
与
\(
\operatorname{vol}(E_{k,t_2}(L))
\)
唯一确定
\[
S(L):=\sum_{i=1}^{k}\log L_i=\log N.
\]
从而在给定 \((p_i)_{i=1}^k\) 的前提下，\(N=\exp(S(L))\) 被两阶体量唯一确定；并且由唯一分解定理，指数向量 \(a\) 在纯算术意义下亦被唯一确定。
\end{corollary}
\begin{proof}
由定义 \(S(L)=\sum_i a_i\log p_i=\log N\)。
另一方面，推论 \ref{cor:xi-godel-loewner-two-level-holographic-inversion} 表明 \(S(L)\) 可由两阶体量作显式反演得到，故 \(\log N\) 亦由两阶体量唯一确定。
在给定素数底 \((p_i)\) 时，唯一分解定理给出从 \(N\) 到指数向量 \(a\) 的唯一性。证毕。
\end{proof}

\begin{proposition}[外幂层间体量的严格幂律：体积塔由一阶体量生成]\label{prop:xi-godel-loewner-volume-renormalization-power-law}
\leanverified{paper\_xi\_godel\_loewner\_volume\_renormalization\_power\_law}
在定理 \ref{thm:xi-godel-loewner-exterior-box-loewner-holographic-weight} 的设定下，记 \(V_{k,t}(L):=\operatorname{vol}(E_{k,t}(L))\)。
则对任意 \(1\le t\le k\) 成立严格幂律关系
\[
V_{k,t}(L)
=
\kappa_{\binom{k}{t}}\binom{k}{t}^{\binom{k}{t}/2}2^{-\binom{k}{t}}
\left(\frac{2^{k}}{\kappa_{k}k^{k/2}}\,V_{k,1}(L)\right)^{\binom{k-1}{t-1}}.
\]
因此固定 \(k\) 时，外幂体量塔 \(\{V_{k,t}(L)\}_{t=1}^{k}\) 由单一量 \(V_{k,1}(L)\) 与仅依赖 \((k,t)\) 的普适常数完全生成。
\end{proposition}
\begin{proof}
由定理 \ref{thm:xi-godel-loewner-exterior-box-loewner-holographic-weight}(2) 在 \(t=1\) 得
\[
V_{k,1}(L)=\kappa_k\,k^{k/2}2^{-k}\prod_{i=1}^{k}L_i,
\qquad
\prod_{i=1}^{k}L_i=\frac{2^{k}}{\kappa_k k^{k/2}}\,V_{k,1}(L).
\]
将其代入同一条体积闭式的一般 \(t\) 情形即得。证毕。
\end{proof}

\begin{theorem}[prime-register 特化下的外幂体量主阶：由圆维支配的二次体量律]\label{thm:xi-godel-prime-register-exterior-volume-quadratic-law}
取素数序列 \((p_i)_{i\ge 1}\)，并对给定 \(k\) 置
\[
L_i:=\log p_i\quad(1\le i\le k),
\qquad
S_k:=\sum_{i=1}^{k}\log\log p_i.
\]
固定 \(t\in\{1,\dots,k\}\)，并令 \(V^{(\mathrm{prime})}_{k,t}:=V_{k,t}(L)\)。
则成立精确恒等式
\[
\log V^{(\mathrm{prime})}_{k,t}
=
\log\kappa_N+\frac{N}{2}\log N-N\log 2+\binom{k-1}{t-1}S_k,
\qquad N=\binom{k}{t}.
\]
并且当 \(t\) 固定、\(k\to\infty\) 时，
\[
\frac{1}{N}\log V^{(\mathrm{prime})}_{k,t}
=
\Bigl(\frac{1}{N}\log\kappa_N+\frac12\log N-\log 2\Bigr)+\frac{t}{k}S_k
=
\Bigl(\tfrac12\log(2\pi e)-\log 2\Bigr)+t\log\log k+O(1),
\]
其中最后一步由素数定理的标准推论 \(S_k=k\log\log k+O(k)\) 给出。
特别地，当 \(t=2\) 时有
\[
\log V^{(\mathrm{prime})}_{k,2}
=
\binom{k}{2}\Bigl(2\log\log k+\tfrac12\log(2\pi e)-\log 2\Bigr)+O(k^2),
\]
从而在圆维语言（此处 \(k=\cdim(\ZZ^k)\)）下呈现二次主阶
\[
\log V^{(\mathrm{prime})}_{\cdim(\ZZ^k),2}
=
\Theta\!\big(\cdim(\ZZ^k)^2\log\log \cdim(\ZZ^k)\big).
\]
\end{theorem}
\begin{proof}
精确恒等式为定理 \ref{thm:xi-godel-loewner-exterior-box-loewner-holographic-weight}(2) 的直接代入，并用
\(
\sum_{I\in\mathcal I_{k,t}}\log L_I=\binom{k-1}{t-1}\sum_{i=1}^k\log L_i
\)
化简得到。
\par
单位维度展开由定理 \ref{thm:xi-godel-loewner-exterior-box-loewner-holographic-weight}(3) 立即给出。
再由 Stirling 公式得到
\(
\frac1N\log\kappa_N+\frac12\log N=\frac12\log(2\pi e)+O(\log N/N)
\)，
并由素数定理推出 \(\log p_k=\log k+O(\log\log k)\)，从而
\(
S_k=k\log\log k+O(k)
\)；
代回即得所示主阶。
最后取 \(t=2\)、\(N=\binom{k}{2}\) 得特别情形。证毕。
\end{proof}

\endinput


\leanverified{paper\_xi\_cdim\_phase\_compression\_aliasing\_exponent}
\begin{theorem}[单位圆相位压缩的混叠指数等于 \texorpdfstring{$r/d$}{r/d}]\label{thm:xi-cdim-phase-compression-aliasing-exponent}
取 \(\TT:=\RR/\ZZ\cong S^1\)，并令
\[
\mathrm{dist}_{\TT^d}(x,y):=\min_{m\in\ZZ^d}|x-y-m|_{\infty}.
\]
对群同态 \(\phi:\ZZ^r\to\TT^d\)，定义
\[
Q_B:=\{u\in\ZZ^r:\ |u|_{\infty}\le B\},
\]
\[
\Delta_B(\phi):=
\min\Bigl\{\mathrm{dist}_{\TT^d}\!\bigl(\phi(u),\phi(v)\bigr):
u\neq v,\ u,v\in Q_B\Bigr\}.
\]
再定义
\[
\alpha_d:=
\inf_{\substack{\phi\in\mathrm{Hom}(\ZZ^r,\TT^d)\\ \phi\ \text{单射}}}
\ \limsup_{B\to\infty}\frac{\log(1/\Delta_B(\phi))}{\log B}.
\]
则
\[
\alpha_d=\frac{r}{d}.
\]
\end{theorem}
\begin{proof}
先证下界。对任意单射 \(\phi\)，\(\phi(Q_B)\) 有 \((2B+1)^r\) 个不同点。
若 \(\Delta_B(\phi)\ge \varepsilon\)，则这些点两两 \(\varepsilon\)-分离。
而 \(\TT^d\) 中任意 \(\varepsilon\)-分离集的基数至多 \(C_d\varepsilon^{-d}\)（由 \(\ell^\infty\)-球装填/网格覆盖）。
故
\[
(2B+1)^r\le C_d\,\Delta_B(\phi)^{-d},
\qquad
\Delta_B(\phi)\le C'_d\,B^{-r/d}.
\]
于是
\[
\limsup_{B\to\infty}\frac{\log(1/\Delta_B(\phi))}{\log B}\ge \frac{r}{d}.
\]
对所有单射 \(\phi\) 取下确界得 \(\alpha_d\ge r/d\)。

再证上界。由经典 Schmidt 型坏可逼近系统存在定理（见 \cite{Schmidt1969BadlyApproximableSystems}），存在矩阵
\[
A\in M_{d\times r}(\RR),\quad c>0,
\]
使得对所有 \(m\in\ZZ^r\setminus\{0\}\) 有
\[
\mathrm{dist}_{\TT^d}(Am,0)\ge c\,|m|_{\infty}^{-r/d}.
\]
定义 \(\phi_A(u):=Au\bmod\ZZ^d\)。上式立即推出 \(\phi_A\) 单射。
对任意不同 \(u,v\in Q_B\)，令 \(m=u-v\neq 0\)，则 \(|m|_{\infty}\le 2B\)，故
\[
\mathrm{dist}_{\TT^d}(\phi_A(u),\phi_A(v))
=
\mathrm{dist}_{\TT^d}(Am,0)
\ge c(2B)^{-r/d}.
\]
即
\[
\Delta_B(\phi_A)\ge c'B^{-r/d}.
\]
于是
\[
\limsup_{B\to\infty}\frac{\log(1/\Delta_B(\phi_A))}{\log B}\le \frac{r}{d},
\]
从而 \(\alpha_d\le r/d\)。
结合上下界得 \(\alpha_d=r/d\)。证毕。
\end{proof}

\begin{corollary}[圆维—精度守恒律：相位压缩的可逆实现下界]\label{cor:xi-cdim-precision-conservation-law}
在定理 \ref{thm:xi-cdim-phase-compression-aliasing-exponent} 的设定下，设实现层相位分辨率为 \(\varepsilon>0\)：当两点距离 \(<\varepsilon\) 时不可区分。
若要求在 \(Q_B\) 上实现 \(\varepsilon\)-意义下可逆解码，则必须满足
\[
\varepsilon\ \lesssim\ B^{-r/d},
\]
等价地，
\[
d\log\frac1\varepsilon\ \gtrsim\ r\log B.
\]
\leanverified{paper_xi_cdim_precision_conservation_law}
\end{corollary}
\begin{proof}
对任意单射 \(\phi\)，由定理 \ref{thm:xi-cdim-phase-compression-aliasing-exponent} 证明中的装填估计，有
\[
\Delta_B(\phi)\le C_d' B^{-r/d}.
\]
若要在分辨率 \(\varepsilon\) 下可逆区分 \(Q_B\) 中所有点，必要条件是 \(\varepsilon<\Delta_B(\phi)\)。
故必有 \(\varepsilon\le C_d'B^{-r/d}\)，即第一式。
两侧取对数并吸收常数得第二式。证毕。
\end{proof}

\begin{remark}[半圆维口径下的同一指数]\label{rem:xi-cdim-plus-phase-precision-index}
若 \(M\) 为有限生成可消去交换幺半群，且其 Grothendieck 群完成 \(G(M)\) 的自由秩为 \(r_M\)，则
\[
\cdim_{+}(M)=\frac{r_M}{2}.
\]
把上述定理应用到 \(G(M)\) 可得对应实现指数
\[
\frac{r_M}{d}=\frac{2\cdim_{+}(M)}{d}.
\]
因此半圆维同样给出可检验的精度收缩率，只是以 \(2\cdim_{+}\) 进入幂指数。
\end{remark}

\endinput

\paragraph{相位实现秩、有限商增长与预算障碍的统一刻画}
沿用定理 \ref{thm:xi-cdim-zeta-euler-abscissa-locality}、
\ref{thm:xi-circle-dimension-continuous-noncompressibility}、
\ref{cor:xi-cdim-precision-conservation-law}、
\ref{thm:xi-godel-integerization-log-overhead}、
\ref{thm:xi-offcritical-endpoint-resolution-lower-bound}、
\ref{cor:xi-comoving-prefix-bit-budget-gain}、
\ref{thm:pom-minkowski-budget-barrier}、
\ref{cor:pom-multi-axis-near1-pole-barrier} 与
\ref{cor:pom-multi-axis-pole-drift-binding} 的记号。

\begin{theorem}[相位实现秩的三重刻画]\label{thm:xi-phase-implementation-rank-triple-characterization}
设 \(G\) 为有限生成离散交换群，并定义
\[
\rho_{\mathrm{ph}}(G):=
\min\Bigl\{
k\in\NN\cup\{0\}:\ \exists\, D\ \text{为零维紧交换群，使}\ 
\widehat G\hookrightarrow \TT^k\times D\ \text{连续单射}
\Bigr\}.
\]
则
\[
\rho_{\mathrm{ph}}(G)
=
\cdim(G)
=
\lim_{n\to\infty}\frac{\log|G/nG|}{\log n}.
\]
\end{theorem}
\begin{proof}
由有限生成交换群结构定理，可写
\[
G\cong \ZZ^r\oplus F,
\]
其中 \(F\) 有限。于是
\[
\cdim(G)=r,
\qquad
|G/nG|=n^r\,|F/nF|.
\]
由于 \(1\le |F/nF|\le |F|\)，可得
\[
\frac{\log|G/nG|}{\log n}
=
r+\frac{\log|F/nF|}{\log n}\longrightarrow r,
\]
从而极限等于 \(\cdim(G)\)。

另一方面，Pontryagin 对偶给出
\[
\widehat G\cong \widehat{\ZZ^r}\times \widehat F\cong \TT^r\times \widehat F,
\]
其中 \(\widehat F\) 为有限群，因而是零维紧交换群。这立刻给出
\(
\rho_{\mathrm{ph}}(G)\le r
\)。

反过来，若存在连续单射
\[
\iota:\widehat G\hookrightarrow \TT^k\times D
\]
其中 \(D\) 为零维紧交换群，则限制到闭子群
\(
\TT^r\times\{0\}\subset \widehat G
\)
得到连续单射
\[
\TT^r\hookrightarrow \TT^k\times D.
\]
由定理 \ref{thm:xi-circle-dimension-continuous-noncompressibility}，必有 \(r\le k\)，故
\(
\rho_{\mathrm{ph}}(G)\ge r
\)。
综上 \(\rho_{\mathrm{ph}}(G)=r=\cdim(G)\)，再与上面的极限公式合并即得结论。
\leanverified{paper\_xi\_phase\_implementation\_rank\_triple\_characterization}
\leanverified{paper\_xi\_phase\_implementation\_rank\_limit\_seeds}
\leanverified{paper\_xi\_phase\_implementation\_rank\_limit\_package}
\end{proof}

\begin{theorem}[有限商标量 Gödel 化的不可约 \texorpdfstring{$\log\log n$}{loglogn} 乘子下界]\label{thm:xi-finite-quotient-scalar-godel-loglog-overhead}
设 \(G\) 为有限生成离散交换群，记
\[
r:=\cdim(G)\ge 1.
\]
固定有限字母表 \(\mathcal A\subset \NN_{\ge 1}\)，并记 \(A:=|\mathcal A|\ge 2\)。
对每个 \(n\ge 2\)，设存在注入编码
\[
\operatorname{enc}_n:G/nG\hookrightarrow \mathcal A^{T_n}.
\]
对任意
\(
a=(a_1,\dots,a_{T_n})\in \mathcal A^{T_n}
\)
定义 Gödel 整数
\[
\Gamma_n(a):=\prod_{t=1}^{T_n}p_t^{a_t},
\]
其中 \(p_t\) 为第 \(t\) 个素数。再记
\[
L_n:=\max_{x\in G/nG}\log_2\Gamma_n(\operatorname{enc}_n(x)).
\]
则当 \(n\to\infty\) 时，
\[
L_n\ \ge\ \frac{r}{\log A}\,\log n\,\log_2\log n\ -\ O(\log n).
\]
\leanverified{paper\_xi\_finite\_quotient\_scalar\_godel\_loglog\_overhead\_seeds}
\end{theorem}
\begin{proof}
仍写
\(
G\cong \ZZ^r\oplus F
\)
其中 \(F\) 有限。于是
\[
|G/nG|=n^r\,|F/nF|,
\qquad
1\le |F/nF|\le |F|.
\]
由编码的注入性，
\[
A^{T_n}\ge |G/nG|,
\]
从而
\[
T_n\log A\ge \log|G/nG|=r\log n+O(1).
\]
故
\[
T_n\ge \frac{r}{\log A}\log n-O(1).
\]

另一方面，对每个码字 \(a\in\mathcal A^{T_n}\subset \NN_{\ge 1}^{T_n}\)，
定理 \ref{thm:xi-godel-integerization-log-overhead} 给出
\[
\log_2\Gamma_n(a)\ge T_n\log_2 T_n-c_0T_n-c_1
\]
（其中 \(c_0,c_1\) 为绝对常数）。因此
\[
L_n\ge T_n\log_2 T_n-c_0T_n-c_1.
\]
再由
\[
T_n=\frac{r}{\log A}\log n+O(1),
\qquad
\log_2 T_n=\log_2\log n+O(1),
\]
便得到
\[
L_n
\ge
\Bigl(\frac{r}{\log A}\log n-O(1)\Bigr)\bigl(\log_2\log n+O(1)\bigr)-O(\log n)
=
\frac{r}{\log A}\log n\,\log_2\log n-O(\log n).
\]
证毕。
\end{proof}

\begin{theorem}[有限商的相位压缩精度下界]\label{thm:xi-finite-quotient-phase-compression-precision-lower-bound}
\leanverified{paper\_xi\_finite\_quotient\_phase\_compression\_precision\_lower\_bound}
设 \(G\) 为有限生成离散交换群，记 \(r:=\cdim(G)\ge 1\)，并写
\[
G\cong \ZZ^r\oplus T
\]
其中 \(T\) 有限。再固定 \(d\ge 1\)，并考虑任意群同态
\[
\phi:\ZZ^r\to \TT^d.
\]
对每个 \(n\ge 2\)，令
\[
B_{n-1}:=\{0,1,\dots,n-1\}^r\subset \ZZ^r.
\]
若要求 \(\phi\) 在 \(B_{n-1}\) 上实现唯一反演，并且每个相位坐标仅保留 \(b_n\) 比特精度，
即总分辨率满足
\[
\delta_n=2^{-b_n}
\]
（按 \(\TT^d\) 的 \(\ell^\infty\) 距离口径），则必有
\[
b_n\ \ge\ \frac{r}{d}\log_2 n-O(1).
\]
\end{theorem}
\begin{proof}
由群同态性，\(\phi\) 在 \(B_{n-1}\) 上可逆等价于：对任意不同的 \(u,v\in B_{n-1}\)，都有
\[
\operatorname{dist}_{\TT^d}\bigl(\phi(u),\phi(v)\bigr)>\delta_n.
\]
而差集满足
\[
B_{n-1}-B_{n-1}=\{-(n-1),\dots,n-1\}^r.
\]
因此上式等价于对任意
\(
0\neq m\in\{-(n-1),\dots,n-1\}^r
\)
都有
\[
\operatorname{dist}_{\TT^d}(\phi(m),0)>\delta_n.
\]
这正是推论 \ref{cor:xi-cdim-precision-conservation-law} 在尺度 \(B=n-1\) 下的可逆实现要求。
于是必有
\[
\delta_n\lesssim (n-1)^{-r/d}.
\]
两边取 \(-\log_2\) 即得
\[
b_n\ge \frac{r}{d}\log_2(n-1)-O(1)
=
\frac{r}{d}\log_2 n-O(1).
\]

最后，由
\[
G/nG\cong (\ZZ/n\ZZ)^r\oplus (T/nT)
\]
可知 \(B_{n-1}\) 恰参数化了 \(G/nG\) 的自由部分代表，故该界正是有限商自由部分相位压缩所必需的精度门槛。
\end{proof}

\begin{theorem}[圆维驱动的 near-\texorpdfstring{$1$}{1} 极点贴边律]\label{thm:xi-cdim-driven-near1-pole-approach-law}
设 \(G\) 为有限生成离散交换群，记
\[
r:=\cdim(G)\ge 1,
\qquad
q_G(n):=|G/nG|.
\]
再固定整数 \(d\ge 1\)。对每个 \(n\)，取一个 \(d\)-轴模数组
\[
m(n)=\bigl(m_1(n),\dots,m_d(n)\bigr)
\]
并记其总预算为
\[
B_n:=\prod_{j=1}^{d}m_j(n).
\]
若 \(B_n\asymp q_G(n)\)，则对充分大的 \(n\)，存在与该预算对应的某个非平凡扭曲 \(\vartheta_n\)，使得
\[
1-\Re s(\vartheta_n)=O\!\bigl(n^{-2r/d}\bigr),
\qquad
|\Delta_{\mathfrak M}(\vartheta_n)|=O\!\bigl(n^{-2r/d}\bigr).
\]
\end{theorem}
\begin{proof}
写
\[
G\cong \ZZ^r\oplus F
\]
其中 \(F\) 有限。于是
\[
q_G(n)=|G/nG|=n^r\,|F/nF|=:n^r u(n),
\qquad
1\le u(n)\le |F|.
\]
故
\[
q_G(n)\asymp n^r.
\]
由假设 \(B_n\asymp q_G(n)\)，立得
\[
B_n^{-2/d}\asymp q_G(n)^{-2/d}\asymp n^{-2r/d}.
\]

另一方面，推论 \ref{cor:pom-multi-axis-near1-pole-barrier} 保证：对每个充分大的预算 \(B_n\)，存在某个非平凡扭曲 \(\vartheta_n\) 使
\[
1-\Re s(\vartheta_n)=O(B_n^{-2/d}).
\]
再由推论 \ref{cor:pom-multi-axis-pole-drift-binding}，同一扭曲满足
\[
|\Delta_{\mathfrak M}(\vartheta_n)|=O(B_n^{-2/d}).
\]
将 \(B_n^{-2/d}\asymp n^{-2r/d}\) 代入即得所示结论。
\leanverified{paper\_xi\_cdim\_driven\_near1\_pole\_approach\_law}
\end{proof}

\begin{theorem}[地址前缀与相位圆因子的正交资源律]\label{thm:xi-prefix-phase-circle-orthogonal-resource-law}
\leanverified{paper\_xi\_prefix\_phase\_circle\_orthogonal\_resource\_law}
固定 \(\delta_0\in(0,\tfrac12)\)，并记
\[
p_{\mathrm{fix}}(T,\delta_0):=
\log_2\!\Bigl(\frac{T^2+(1+\delta_0)^2}{4\delta_0}\Bigr),
\qquad
p_{\mathrm{com}}(\delta_0):=
\log_2\!\Bigl(\frac{(1+\delta_0)^2}{4\delta_0}\Bigr).
\]
则当 \(T\to\infty\) 时，
\[
p_{\mathrm{fix}}(T,\delta_0)-p_{\mathrm{com}}(\delta_0)=2\log_2T+O(1).
\]

另一方面，设 \(K\) 为紧 Hausdorff 空间，并假设存在连续单射
\[
j:\TT^r\hookrightarrow K.
\]
若 \(D\) 为任意零维紧交换群，且
\[
\Psi:K\hookrightarrow \TT^k\times D
\]
为连续单射，则必有
\[
k\ge r.
\]
因此，共动前缀、\(\mathsf{NULL}\) 证书以及任意零维账本至多回收端点位预算，而不能减少连续相位圆因子的最小数目。
\end{theorem}
\begin{proof}
由定理 \ref{thm:xi-offcritical-endpoint-resolution-lower-bound} 与推论 \ref{cor:xi-comoving-prefix-bit-budget-gain}，
\(
p_{\mathrm{fix}}(T,\delta_0)
\)
与
\(
p_{\mathrm{com}}(\delta_0)
\)
分别正是固定图表与共动前缀图表下的必要位预算下界，故
\[
p_{\mathrm{fix}}(T,\delta_0)-p_{\mathrm{com}}(\delta_0)
=
\log_2\!\Bigl(1+\frac{T^2}{(1+\delta_0)^2}\Bigr)
=
2\log_2 T+O(1).
\]

再看相位圆因子数。若 \(\Psi\) 存在，则复合映射
\[
\TT^r\xrightarrow{\ j\ } K \xrightarrow{\ \Psi\ } \TT^k\times D
\]
仍为连续单射。由定理 \ref{thm:xi-circle-dimension-continuous-noncompressibility}，
任意零维紧因子 \(D\) 都不能替代连续圆方向，故必有 \(r\le k\)。

于是，前缀外置所作用的是端点/地址层的离散预算，而连续相位圆因子的最小数目仍由 \(r\) 强制。
\(\mathsf{NULL}\) 证书与其它离散账本都可并入零维因子 \(D\)，因此不会改变该下界。
\end{proof}

\endinput

在圆维签名的口径下，有限 torsion 账本既可作为实现层的外置寄存器，也可由内禀算术数据以规范方式生成。
以下给出一条以判别群（cokernel）为核心的构造链：对称整双线性型的判别群与其对偶紧环面同态的覆盖核自然同构；对数域迹格，该判别群等同于 codifferent 商并由判别式决定坏素支撑；局部 Smith 正规形进一步给出按素数分解的“轴数—深度”离散谱；最后，判别账本携带的 \(\QQ/\ZZ\)-值配对诱导规范有限傅里叶算子，并与素分解严格张量化兼容。

\begin{theorem}[判别群与覆盖核：最小有限账本的规范实现]\label{thm:xi-disc-ledger-kernel-dual}
\leanverified{paper\_xi\_disc\_ledger\_kernel\_dual}
设 \(M\) 为秩 \(n\) 的自由 \(\ZZ\)-模，\(B:M\times M\to \ZZ\) 为对称双线性型，且其 \(\QQ\)-线性延拓在 \(M_{\QQ}:=M\otimes_{\ZZ}\QQ\) 上非退化。
定义同态
\[
\varphi_B:M\to M^\vee:=\operatorname{Hom}_{\ZZ}(M,\ZZ),\qquad
x\mapsto (y\mapsto B(x,y)),
\]
并定义判别群
\[
D_B:=\operatorname{coker}(\varphi_B)=M^\vee/\varphi_B(M).
\]
令紧环面
\[
\TT_M:=\operatorname{Hom}(M,\RR/\ZZ)\cong (\RR/\ZZ)^n.
\]
则：
\begin{enumerate}
  \item \(\varphi_B\) 在 \(\QQ\)-上线性同构，故 \(\widehat{\varphi_B}:\TT_M\to \TT_M\) 为有限覆盖（isogeny），并存在自然同构
  \[
  \ker(\widehat{\varphi_B})\ \cong\ \widehat{D_B}:=\operatorname{Hom}(D_B,\RR/\ZZ),
  \]
  从而 \(|\ker(\widehat{\varphi_B})|=|D_B|\)。
  \item 设 \(R\) 为有限生成离散交换群且 \(r:\TT_M\to \widehat R\) 为连续群同态，使得
  \[
  (\widehat{\varphi_B},r):\TT_M\longrightarrow \TT_M\times \widehat R
  \]
  为单射，则 \(\widehat R\) 必含一个与 \(\ker(\widehat{\varphi_B})\) 同构的有限子群；等价地，\(R\) 的 torsion 商中必出现与 \(D_B\) 同阶的因子。特别地，任何此类实现的不可规避素数支撑包含 \(|D_B|\) 的全部素因子（等价包含 \(|\det(B)|\) 的全部素因子）。
\end{enumerate}
\end{theorem}
\begin{proof}
取 \(\ZZ\)-基使 \(M\cong \ZZ^n\)。此时 \(\varphi_B\) 由整矩阵 \(A\in M_n(\ZZ)\) 表示，且
\(
D_B\cong \ZZ^n/A\ZZ^n
\)。
对短正合列
\[
0\to \ZZ^n \xrightarrow{A} \ZZ^n \to \ZZ^n/A\ZZ^n\to 0
\]
取 Pontryagin 对偶（离散群的对偶函子保持正合），得
\[
0\to \widehat{\ZZ^n/A\ZZ^n}\to \widehat{\ZZ^n}\xrightarrow{\widehat A}\widehat{\ZZ^n}\to 0.
\]
又 \(\widehat{\ZZ^n}\cong (\RR/\ZZ)^n\cong \TT_M\) 且 \(\widehat A\) 即由 \(A^{\mathsf T}\) 在 \((\RR/\ZZ)^n\) 上诱导的覆盖同态；由于 \(B\) 对称，可在更换基后取 \(A=A^{\mathsf T}\)，从而得到 (1)。
\par
对 (2)，若 \((\widehat{\varphi_B},r)\) 单射，则 \(r\) 在 \(\ker(\widehat{\varphi_B})\) 上亦为单射，故 \(\ker(\widehat{\varphi_B})\) 可嵌入 \(\widehat R\)。
而任意紧交换群 \(\widehat R\) 的有限子群对偶于 \(R\) 的 torsion 商，因此 \(R\) 的 torsion 部分必须包含与 \(\ker(\widehat{\varphi_B})\) 同构的对偶信息；由 (1) 得其阶为 \(|D_B|\)，素数支撑断言随之成立。
\end{proof}

\begin{theorem}[迹格的判别账本与坏素支撑]\label{thm:xi-trace-lattice-discriminant-ledger}
令 \(K/\QQ\) 为次数 \(n\) 的数域，\(\mathcal O_K\) 为其整数环。
取迹配对
\[
B(x,y):=\operatorname{Tr}_{K/\QQ}(xy),\qquad x,y\in\mathcal O_K,
\]
并定义迹对偶格（codifferent）
\[
\mathcal O_K^\vee:=\{\alpha\in K:\operatorname{Tr}_{K/\QQ}(\alpha\,\mathcal O_K)\subseteq \ZZ\}.
\]
令
\[
D_K:=\mathcal O_K^\vee/\mathcal O_K.
\]
则存在自然同构 \(D_K\cong D_B\)，并且
\[
|D_K|=|\Delta_K|,
\]
其中 \(\Delta_K\) 为 \(K\) 的绝对判别式。
从而 \(\widehat{\varphi_B}:\TT_{\mathcal O_K}\to\TT_{\mathcal O_K}\) 的覆盖核同构于 \(\widehat{D_K}\)，且任何使该同态在扩展保存口径下可逆的审计实现都不可避免地引入坏素支撑
\[
\{p:\ p\mid \Delta_K\}.
\]
\end{theorem}
\begin{proof}
对任意整基 \((e_i)\) 的迹矩阵 \(P=(\operatorname{Tr}(e_ie_j))\) 有 \(\det(P)=\Delta_K\)。
而 \([\mathcal O_K^\vee:\mathcal O_K]=|\det(P)|\) 为经典格论恒等式，故 \(|D_K|=[\mathcal O_K^\vee:\mathcal O_K]=|\Delta_K|\)。
同构 \(D_K\cong D_B\) 来自 \(\varphi_B\) 与对偶格的定义对齐。
最后的坏素支撑结论由定理 \ref{thm:xi-disc-ledger-kernel-dual}(2) 立即推出。
\end{proof}

\begin{proposition}[局部最小 \(p\)-进轴数与迹矩阵秩亏]\label{prop:xi-disc-ledger-local-axis-rank-defect}
沿用定理 \ref{thm:xi-trace-lattice-discriminant-ledger} 的记号，取 \(\mathcal O_K\) 的任意整基并令 \(P=(\operatorname{Tr}(e_ie_j))\in M_n(\ZZ)\)。
对每个素数 \(p\) 定义
\[
g_p(K):=\dim_{\FF_p}\bigl(D_K/pD_K\bigr).
\]
则有显式公式
\[
g_p(K)=n-\operatorname{rank}_{\FF_p}(P\bmod p).
\]
并且 \(g_p(K)\) 等于 \(D_K\) 的 \(p\)-主分量所需的最小生成元个数，从而给出任何审计实现中不可规避的独立 \(p\)-进外置轴数量下界。
\end{proposition}
\begin{proof}
由 \(D_K\cong \ZZ^n/P\ZZ^n\) 得
\[
D_K/pD_K\cong (\ZZ^n/P\ZZ^n)\otimes_{\ZZ}\FF_p\cong \FF_p^n/(P\bmod p)\FF_p^n,
\]
故其维数为 \(n-\operatorname{rank}_{\FF_p}(P\bmod p)\)。
有限阿贝尔群 \(G\) 的 \(p\)-秩 \(r_p(G):=\dim_{\FF_p}(G/pG)\) 等于其 \(p\)-主分量的最小生成元个数，代入 \(G=D_K\) 即得结论。
\end{proof}

\begin{theorem}[判别素账本的轴数—深度全分解]\label{thm:xi-disc-ledger-axis-depth-smith}
\leanverified{paper\_xi\_disc\_ledger\_axis\_depth\_smith}
沿用上述记号。
对每个素数 \(p\)，令 \(P\) 在 \(\ZZ_p\) 上的 Smith 正规形为
\[
UPV=\operatorname{diag}(p^{\alpha_{p,1}},\dots,p^{\alpha_{p,n}}),
\qquad
U,V\in \operatorname{GL}_n(\ZZ_p),
\qquad
0\le \alpha_{p,n}\le\cdots\le \alpha_{p,1},
\]
并令
\[
g_p(K):=\#\{i:\alpha_{p,i}>0\},\qquad
\boldsymbol{\alpha}_p(K):=(\alpha_{p,1}\ge\cdots\ge\alpha_{p,g_p(K)}\ge 1).
\]
则：
\begin{enumerate}
  \item \((D_K)_{(p)}\cong \bigoplus_{i=1}^{g_p(K)}\ZZ/p^{\alpha_{p,i}}\ZZ\)，并且序列 \(\boldsymbol{\alpha}_p(K)\) 与整基选择无关；
  \item 判别式的 \(p\)-进指数满足恒等式
  \[
  v_p(\Delta_K)=\sum_{i=1}^{g_p(K)}\alpha_{p,i}.
  \]
\end{enumerate}
因此 \(g_p(K)\) 给出不可规避的 \(p\)-进外置轴数目，而 \(\boldsymbol{\alpha}_p(K)\) 给出每一轴的必需深度；二者共同构成判别素账本的完备离散谱。
\end{theorem}
\begin{proof}
Smith 正规形在 PID \(\ZZ_p\) 上存在且不变量唯一，给出 cokernel 的直和分解与基不变性。
另一方面，对角元之积等于 \(\det(P)\)，而 \(\det(P)=\Delta_K\)，故
\(
v_p(\Delta_K)=\sum_{i=1}^{n}\alpha_{p,i}
=\sum_{i=1}^{g_p(K)}\alpha_{p,i}
\)。
\end{proof}

\begin{corollary}[二次数域的循环判别账本]\label{cor:xi-disc-ledger-quadratic-cyclic}
若 \([K:\QQ]=2\)，则对每个坏素 \(p\mid \Delta_K\) 有 \(g_p(K)=1\) 且
\[
(D_K)_{(p)}\cong \ZZ/p^{v_p(\Delta_K)}\ZZ.
\]
从而 \(D_K\cong \ZZ/|\Delta_K|\ZZ\) 为循环群；审计意义下每个坏素仅需单一 \(p\)-进外置轴，其深度恰为 \(v_p(\Delta_K)\)。
\end{corollary}
\begin{proof}
二次情形下 \(P\in M_2(\ZZ)\) 且 \(\det(P)\neq 0\)。其 Smith 不变量在同构意义下为 \((1,|\det(P)|)\)，故 cokernel 循环；局部化到 \(\ZZ_p\) 后仅出现一个非零指数 \(v_p(\det(P))=v_p(\Delta_K)\)。
\end{proof}

\begin{proposition}[判别账本的规范有限傅里叶算子与素分解张量化]\label{prop:xi-disc-ledger-fourier-tensorization}
令 \(D\) 为有限阿贝尔群，\(\langle\cdot,\cdot\rangle:D\times D\to \QQ/\ZZ\) 为非退化双线性配对，并定义双角色
\[
\chi(a,b):=\exp\!\bigl(2\pi \mathrm{i}\,\langle a,b\rangle\bigr).
\]
在群代数 \(\CC[D]\) 上定义算子
\[
(\mathcal F_\chi f)(a):=|D|^{-1/2}\sum_{b\in D}\chi(a,b)\,f(b).
\]
则：
\begin{enumerate}
  \item \(\mathcal F_\chi\) 为酉算子，且仅由配对 \((D,\langle\cdot,\cdot\rangle)\) 决定；
  \item 若 \(D=\bigoplus_{p}D_{(p)}\) 为 \(p\)-主分解且配对按该直和正交分解，则存在规范张量同构
  \[
  \CC[D]\cong \bigotimes_p \CC[D_{(p)}],
  \qquad
  \mathcal F_\chi\cong \bigotimes_p \mathcal F_{\chi_{(p)}},
  \]
  其中 \(\chi_{(p)}\) 为 \(\chi\) 在 \(D_{(p)}\) 上的限制。
\end{enumerate}
在定理 \ref{thm:xi-trace-lattice-discriminant-ledger} 的设定下，判别账本 \(D_K\) 上存在由迹对偶格诱导的自然 \(\QQ/\ZZ\)-值配对，从而 \(\mathcal F_{\chi_K}\) 的谱结构严格张量化分解为各坏素的局部谱因子。
\end{proposition}
\begin{proof}
对任意 \(a,a'\in D\)，利用非退化性可得字符正交关系
\[
\sum_{b\in D}\chi(a,b)\overline{\chi(a',b)}=
\sum_{b\in D}\chi(a-a',b)=
\begin{cases}
|D|,& a=a',\\
0,& a\neq a',
\end{cases}
\]
从而矩阵 \((|D|^{-1/2}\chi(a,b))\) 为酉矩阵，即得 (1)。
对 (2)，当 \(D=\bigoplus_p D_{(p)}\) 且配对正交分解时，双角色按素因子化：\(\chi=\prod_p \chi_{(p)}\)。
将求和按直积拆分并用 \(|D|^{1/2}=\prod_p |D_{(p)}|^{1/2}\) 得到算子的张量分解，群代数的张量同构由有限直和的标准基给出。
最后一段为定义对齐与前述分解的直接应用。
\end{proof}

\begin{conjecture}[判别账本相位谱的四点均匀性]\label{conj:xi-disc-ledger-fourier-spectrum-fourpoint}
在命题 \ref{prop:xi-disc-ledger-fourier-tensorization} 的设定下，进一步假设配对 \(\langle\cdot,\cdot\rangle\) 对称。
则 \(\mathcal F_\chi^4=\mathrm{Id}\)，从而 \(\mathcal F_\chi\) 的谱包含于 \(\{\pm 1,\pm \mathrm{i}\}\)。
固定次数 \(n\ge 2\)，令 \(K\) 在所有 \([K:\QQ]=n\) 且 \(|\Delta_K|\to\infty\) 的数域族中变化，并令 \(\mu_K\) 为 \(\mathcal F_{\chi_K}\) 的经验谱测度（按特征值计重）。
猜想 \(\mu_K\) 在弱-\(\ast\) 意义下趋于 \(\{\pm 1,\pm \mathrm{i}\}\) 上的均匀测度。
\end{conjecture}


\begin{theorem}[最大纤维生成乘法半群的圆维无界]\label{thm:xi-max-fiber-monoid-circle-dimension-unbounded}
在折叠语义下令 \(D_m:=\max_{x\in X_m}d_m(x)\) 为分辨率 \(m\) 的最大纤维多重度。
对 \(N\ge 1\) 定义乘法幺半群
\[
\mathcal{M}_N:=\langle D_1,D_2,\dots,D_N\rangle_{\times}\ \subset\ \NN_{\ge 1},
\]
并令 \(G_N\subset \QQ_{>0}^{\times}\) 为其群完成（等价地：由 \(\mathcal{M}_N\) 生成的乘法子群）。
则 \(\cdim(G_N)\) 随 \(N\) 取值无界。
\end{theorem}
\begin{proof}
对任意有限集合 \(S\subset\NN_{\ge 1}\)，由素数赋值分解得 \(\langle S\rangle\subset\QQ_{>0}^{\times}\) 为有限生成自由交换群，
其自由秩等于出现于 \(\{s:\ s\in S\}\) 的不同素数个数；由有限生成交换群的圆维公式 \(\cdim(G)=\dim_{\QQ}(G\otimes_{\ZZ}\QQ)\)，该自由秩即等于 \(\cdim(\langle S\rangle)\)。
因此若记
\(
\mathcal{P}_N:=\{p\ \text{素数}:\ p\mid D_m\ \text{对某个}\ m\le N\}
\)，
则 \(\cdim(G_N)=|\mathcal{P}_N|\)。
\par
由最大纤维的闭式（定理 \ref{thm:pom-max-fiber}），对 \(k\ge 0\) 有
\(
D_{2k}=F_{k+2}
\) 与
\(
D_{2k+1}=2F_{k+1}
\)。
而 Fibonacci 数列满足原始素因子强迫：除有限例外外，每个 \(F_n\) 存在一素数 \(p_n\mid F_n\) 且 \(p_n\nmid \prod_{1\le j<n}F_j\)（Zsigmondy/Carmichael 型结论；亦见命题 \ref{prop:pom-collision-aq-fibonacci-bf-primitive-prime-divisor} 的同类用法）。
取 \(n=k+2\) 得：当 \(k\) 充分大时，\(D_{2k}=F_{k+2}\) 含有一个此前未出现的素数因子。
因此 \(|\mathcal{P}_{2k}|\to\infty\)，从而 \(\cdim(G_N)=|\mathcal{P}_N|\) 随 \(N\) 无界。证毕。
\end{proof}

\paragraph{有限原子零色散：常数项曲线的闭式与折点结构}
在零色散口径下，常数偏置尺度的最优成功率极限可由一个正随机变量 \(Z\) 统一刻画：
令
\begin{equation}\label{eq:xi-zero-dispersion-psi-def}
\Psi(c):=\EE\!\left[\min\!\left\{1,\frac{e^{c}}{Z}\right\}\right]\qquad(c\in\RR),
\end{equation}
并定义常数项函数为右逆阈值
\begin{equation}\label{eq:xi-zero-dispersion-b-eps-def}
b(\varepsilon):=\inf\left\{c\in\RR:\ \Psi(c)\ge 1-\varepsilon\right\}\qquad(\varepsilon\in(0,1)).
\end{equation}
当 \(Z\) 为有限原子分布时，\(\Psi\) 的分段仿射结构使 \(b(\varepsilon)\) 可显式反演。

\begin{theorem}[有限原子零色散下的常数项闭式：\texorpdfstring{$b(\varepsilon)$}{b(eps)} 的分段反演]\label{thm:xi-zero-dispersion-finite-atom-b-eps}
\leanverified{paper\_xi\_zero\_dispersion\_finite\_atom\_b\_eps}
设 \(Z\) 为正随机变量且满足 \(\EE[Z^{-1}]<\infty\)。
若 \(Z\) 为有限原子分布：存在 \(\kappa\ge 1\)、权重 \(p_i>0\) 与原子位置 \(0<z_1<\cdots<z_\kappa\) 使
\[
\PP(Z=z_i)=p_i,\qquad \sum_{i=1}^{\kappa}p_i=1,
\]
则令 \(s:=1-\varepsilon\) 并定义累积质量与加权尾和
\[
P_j:=\sum_{i=1}^{j}p_i\quad(P_0:=0),\qquad
Q_j:=\sum_{i=j+1}^{\kappa}\frac{p_i}{z_i}\quad(Q_\kappa:=0).
\]
则由 \eqref{eq:xi-zero-dispersion-b-eps-def} 定义的 \(b(\varepsilon)\) 具有分段闭式表达：对每个 \(j\in\{0,1,\dots,\kappa-1\}\)，当
\[
P_j<s\le P_{j+1}
\]
时，有
\[
b(\varepsilon)=\log\!\left(\frac{s-P_j}{Q_j}\right).
\]
并且将 \(b\) 以连续方式延拓到 \(\varepsilon=0\) 时满足
\(
b(0)=\log z_\kappa
\)。
\end{theorem}
\begin{proof}
由有限原子表示与 \eqref{eq:xi-zero-dispersion-psi-def}，
\[
\Psi(c)=\sum_{i=1}^{\kappa}p_i\,\min\!\left\{1,\frac{e^c}{z_i}\right\}.
\]
对任意 \(c\)，设 \(j=j(c)\) 满足 \(z_j\le e^c<z_{j+1}\)（约定 \(z_0:=0\)、\(z_{\kappa+1}:=+\infty\)）。
则当 \(i\le j\) 时 \(\min\{1,e^c/z_i\}=1\)，当 \(i>j\) 时 \(\min\{1,e^c/z_i\}=e^c/z_i\)。
从而
\[
\Psi(c)=\sum_{i\le j}p_i+e^c\sum_{i>j}\frac{p_i}{z_i}=P_j+e^cQ_j.
\]
在区间 \(c\in[\log z_j,\log z_{j+1})\) 上，\(\Psi\) 关于 \(e^c\) 为仿射函数，求解 \(\Psi(c)=s\) 得
\(
e^c=(s-P_j)/Q_j
\)，即所示闭式。
由 \(\Psi\) 单调性，解即为 \(\inf\{c:\Psi(c)\ge s\}\)。
当 \(s=1\) 时需 \(\Psi(c)=1\)，在有限原子情形等价于 \(e^c\ge z_\kappa\)，故最小为 \(c=\log z_\kappa\)。证毕。
\end{proof}

\begin{corollary}[折点结构：\texorpdfstring{$\Psi$}{Psi} 的折点与 \texorpdfstring{$b(\varepsilon)$}{b(eps)} 的不可微点]\label{cor:xi-zero-dispersion-kinks}
在定理 \ref{thm:xi-zero-dispersion-finite-atom-b-eps} 的设定下，令 \(t:=e^c\) 并记 \(\Phi(t):=\Psi(\log t)\)。
则：
\begin{enumerate}
  \item \(\Phi\) 关于 \(t\) 为凹的分段线性函数，线性片段数为 \(\kappa\)，折点恰位于 \(t=z_1,\dots,z_\kappa\)（等价 \(c=\log z_1,\dots,\log z_\kappa\)）。
  \item \(b(\varepsilon)\) 在 \((0,1)\) 上连续、严格递增且分段光滑；其不可微点恰位于
  \[
  \varepsilon=1-P_1,\ 1-P_2,\ \dots,\ 1-P_{\kappa-1}.
  \]
\end{enumerate}
\end{corollary}
\begin{proof}
由证明中 \(\Psi(c)=P_j+e^cQ_j\) 的分段仿射表示，折点发生在 \(e^c=z_i\) 处，得 (1)。
由 \(\Psi\) 连续严格递增可知 \(b\) 为其（右）逆阈值映射 \(b(\varepsilon)=\inf\{c:\Psi(c)\ge 1-\varepsilon\}\)，从而连续严格递增并在每个仿射片段上光滑；不可微点对应于 \(\Psi\) 的折点处之逆像，即 \(s=P_j\) 等价于 \(\varepsilon=1-P_j\)。证毕。
\end{proof}

\begin{theorem}[零色散侧信息曲线的可识别性：由 \texorpdfstring{$b(\varepsilon)$}{b(eps)} 重构极限律]\label{thm:xi-zero-dispersion-identifiability}
设 \(Z\) 为正随机变量并满足 \(\EE[Z^{-1}]<\infty\)。
定义 \(\Psi\) 如 \eqref{eq:xi-zero-dispersion-psi-def}，并令 \(\Phi(t):=\Psi(\log t)=\EE[\min\{1,t/Z\}]\)（\(t>0\)）。
则：
\begin{enumerate}
  \item \(\Phi\) 在 \((0,\infty)\) 上连续、单调非减且关于 \(t\) 为凹函数；\(\Psi\) 在 \(\RR\) 上连续、单调非减且关于 \(e^c\) 为凹函数。
  \item 知道 \(\Psi\)（等价地知道 \(b(\varepsilon)\) 在 \((0,1)\) 上的全部取值）即可唯一确定 \(Z\) 的分布。
        更精确地，定义下降函数
        \begin{equation}\label{eq:xi-zero-dispersion-phi-derivative}
        q(t):=\EE\!\left[\frac{1}{Z}\,\ind_{\{Z>t\}}\right]\qquad(t>0).
        \end{equation}
        则 \(q\) 单调不增且右连续，并且 \(\Phi\) 在几乎处处可导且满足 \(\Phi'(t)=q(t)\)。
        令 \(q\) 的 Stieltjes 测度为 \(-\dd q\)，则 \(Z\) 的分布函数 \(F_Z\) 满足测度恒等式
        \begin{equation}\label{eq:xi-zero-dispersion-measure-inversion}
        \dd F_Z(t)=t\,(-\dd q(t))\qquad\text{于 }(0,\infty).
        \end{equation}
\end{enumerate}
\leanverified{paper\_xi\_zero\_dispersion\_identifiability}
\end{theorem}
\begin{proof}
(1) 对固定样本 \(Z=z>0\)，函数 \(t\mapsto \min\{1,t/z\}\) 连续、单调非减且关于 \(t\) 凹。
取期望保持上述性质，得 \(\Phi\) 连续、单调与凹；从而 \(\Psi(c)=\Phi(e^c)\) 连续单调，且关于 \(e^c\) 凹。
\par
(2) 对固定 \(z>0\)，函数 \(t\mapsto \min\{1,t/z\}\) 在 \(t\neq z\) 处可导且导数为 \(z^{-1}\ind_{\{t<z\}}\)。
由支配收敛（\(\EE[Z^{-1}]<\infty\)）可知 \(\Phi\) 在几乎处处可导且 \(\Phi'(t)=q(t)\)，其中 \(q\) 由 \eqref{eq:xi-zero-dispersion-phi-derivative} 定义。
令 \(\nu\) 为有限测度 \(\nu(\dd z):=z^{-1}\dd F_Z(z)\)，则对一切 \(t>0\) 有
\[
q(t)=\int_{(t,\infty)}\frac{1}{z}\,\dd F_Z(z)=\nu((t,\infty)).
\]
因此 \(\nu\) 由 \(q\) 唯一确定为其 Stieltjes 测度 \(\nu=-\dd q\)。
再由 \(\dd F_Z(t)=t\,\dd\nu(t)\) 得到 \eqref{eq:xi-zero-dispersion-measure-inversion}，从而 \(\Phi\)（等价 \(\Psi\)）唯一确定 \(F_Z\)。
最后，\(b(\varepsilon)\) 是 \(\Psi\) 的右逆阈值映射 \eqref{eq:xi-zero-dispersion-b-eps-def}，故二者互相决定。证毕。
\end{proof}

\begin{corollary}[有限原子情形的显式重构：由折点恢复 \texorpdfstring{$(z_i,p_i)$}{(zi,pi)}]\label{cor:xi-zero-dispersion-finite-atom-reconstruction}
若定理 \ref{thm:xi-zero-dispersion-finite-atom-b-eps} 的 \(Z\) 为有限原子分布，则仅由曲线 \(\Psi\)（等价地由 \(b(\varepsilon)\)）可显式重构全部原子 \(\{(z_i,p_i)\}_{i=1}^{\kappa}\)：
\begin{enumerate}
  \item 原子位置由折点给出：\(z_i=\exp(c_i)\)，其中 \(\{c_i\}\) 为 \(\Psi\) 的全部折点（等价为 \(b\) 的全部不可微点经单调变换回去）。
  \item 在区间 \(c\in(\log z_j,\log z_{j+1})\) 上有 \(\Psi(c)=P_j+e^cQ_j\)，故关于 \(t=e^c\) 的斜率为 \(Q_j\)。
        跨越折点 \(t=z_{j+1}\) 时，尾和下降量满足 \(Q_j-Q_{j+1}=p_{j+1}/z_{j+1}\)，从而递推得到全部 \(p_i\)。
\end{enumerate}
\end{corollary}
\begin{proof}
折点结构由推论 \ref{cor:xi-zero-dispersion-kinks} 给出。
而证明中 \(Q_j\) 的分段常数性与折点处的跳变公式直接导出 \(p_{j+1}/z_{j+1}\)，由 \(p_i\) 归一化即可递推恢复全部权重。证毕。
\end{proof}

\paragraph{串接律：二阶常数项的极小卷积上界}
当审计链条由多个可串接读出步骤组成时，侧信息预算的二阶常数项满足可组合的不等式上界。

\begin{theorem}[二阶常数项的串接上界：极小卷积律]\label{thm:xi-zero-dispersion-infimal-convolution}
考虑两段可串接的读出步骤（可理解为两段审计可追溯读出链）。
对第 \(r\in\{1,2\}\) 段，设在平均错误率 \(\varepsilon\) 下的最小侧信息对数预算满足零色散展开
\[
L^{(r)}_m(\varepsilon)=m\gamma_r+b_r(\varepsilon)+o(1).
\]
则对串接协议，存在 \(\gamma_{\mathrm{comp}},b_{\mathrm{comp}}\) 使
\[
L^{(\mathrm{comp})}_m(\varepsilon)=m\gamma_{\mathrm{comp}}+b_{\mathrm{comp}}(\varepsilon)+o(1),
\]
并且满足不等式
\[
\gamma_{\mathrm{comp}}\le \gamma_1+\gamma_2,
\qquad
b_{\mathrm{comp}}(\varepsilon)\le
\inf_{\varepsilon_1+\varepsilon_2\le \varepsilon}\Bigl[b_1(\varepsilon_1)+b_2(\varepsilon_2)\Bigr].
\]
\leanverified{compositeRate\_3\_5}
\leanverified{compositeRate\_7\_2}
\leanverified{compositeRate\_comm}
\leanverified{compositeRate\_assoc}
\leanverified{splitCost\_2\_3\_4\_6}
\leanverified{splitCost\_equal}
\leanverified{splitCost\_equal\_valid}
\leanverified{totalBudget\_add\_m}
\leanverified{totalBudget\_unit\_diff}
\leanverified{rate\_sub\_additivity}
\leanverified{linear\_min\_seed}
\leanverified{infconv\_at\_zero}
\leanverified{paper\_xi\_zero\_dispersion\_infimal\_convolution\_seeds}
\leanverified{paper\_xi\_zero\_dispersion\_infimal\_convolution\_package}
\end{theorem}
\begin{proof}
取任意 \(\varepsilon_1,\varepsilon_2\ge 0\) 且 \(\varepsilon_1+\varepsilon_2\le\varepsilon\)。
在第 \(m\) 层分别为两段选择预算
\(
M^{(r)}_m=\exp(L^{(r)}_m(\varepsilon_r))
\)。
用相互独立的寄存器同时携带两段侧信息，总预算为 \(M_m:=M^{(1)}_mM^{(2)}_m\)，故
\[
\log M_m=L^{(1)}_m(\varepsilon_1)+L^{(2)}_m(\varepsilon_2)
=m(\gamma_1+\gamma_2)+b_1(\varepsilon_1)+b_2(\varepsilon_2)+o(1).
\]
若将整体失败事件定义为任一段失败，则并集界给出总失败率 \(\le\varepsilon_1+\varepsilon_2\le\varepsilon\)。
因此存在一方案在失败率 \(\le\varepsilon\) 下实现上述预算；对极小预算取下确界并对所有可行分解取下确界，即得所示上界。
主指数不等式为同一预算等式中 \(m\)-系数项的比较。证毕。
\end{proof}

\begin{corollary}[有限原子零色散在串接下的折点数上界]\label{cor:xi-zero-dispersion-kink-count-under-composition}
若两段协议各自对应有限原子极限 \(Z^{(1)},Z^{(2)}\)，其原子数分别为 \(\kappa_1,\kappa_2\)，并记串接协议的常数项函数为 \(b_{\mathrm{comp}}(\varepsilon)\)。
则 \(b_{\mathrm{comp}}\) 的不可微点数满足
\[
\#\mathrm{kink}\bigl(b_{\mathrm{comp}}\bigr)\le (\kappa_1-1)+(\kappa_2-1),
\]
并且在一般位置下可达到该上界。
\end{corollary}
\begin{proof}
由定理 \ref{thm:xi-zero-dispersion-infimal-convolution}，\(b_{\mathrm{comp}}\) 受两侧常数项的极小卷积上界控制。
极小卷积的不可微点只能来源于两侧函数不可微点的对齐，故不可微点数至多为两者之和。证毕。
\end{proof}

\paragraph{二阶圆维签名：把连续圆因子与离散侧信息预算合并为可辨识不变量}
圆维刻意忽略完全不连通的离散寄存器（定义 \ref{def:circle-dimension}），而零色散预算常数项 \(b(\varepsilon)\) 与一阶指数 \(\gamma\) 正是对该“被忽略信息”的定量刻画。
据此引入二阶圆维签名作为可核验的二阶审计复杂度坐标。

\begin{definition}[二阶圆维签名作为侧信息复杂度坐标]\label{def:xi-second-order-circle-dimension-signature}
固定单圆分辨率标尺 \(\lambda>1\)。
设最小侧信息预算满足零色散展开 \(L_m(\varepsilon)=m\gamma+b(\varepsilon)+o(1)\)。
定义\textbf{二阶圆维签名}为
\[
\Sigma^{(2)}(\varepsilon):=\bigl(\cdim,\,\gamma/\log\lambda,\,b(\varepsilon)\bigr),
\]
其中 \(\cdim\) 为结构圆维（定义 \ref{def:circle-dimension}）。
在具体折叠系统中可取 \(\lambda=\varphi\) 作为单圆分辨率标尺。
\end{definition}

\begin{proposition}[有限原子零色散下的完备性：\texorpdfstring{$\Sigma^{(2)}(\varepsilon)$}{Sigma2(eps)} 等价于极限律]\label{prop:xi-second-order-circle-dimension-signature-completeness}
\leanverified{paper\_xi\_second\_order\_circle\_dimension\_signature\_completeness}
在定理 \ref{thm:xi-zero-dispersion-finite-atom-b-eps} 的有限原子零色散条件下，
\(\{\Sigma^{(2)}(\varepsilon):\varepsilon\in(0,1)\}\) 与极限律 \(\mathrm{Law}(Z)\) 等价；
亦即二阶常数项曲线 \(\varepsilon\mapsto b(\varepsilon)\)（配合一阶指数 \(\gamma\) 与结构圆维 \(\cdim\)）给出对“二阶审计复杂度”的完备统计不变量。
\end{proposition}
\begin{proof}
由定理 \ref{thm:xi-zero-dispersion-identifiability}，\(\{b(\varepsilon)\}_{\varepsilon\in(0,1)}\) 与 \(\Psi\) 互相决定，且 \(\Psi\) 唯一确定 \(\mathrm{Law}(Z)\)。
一阶指数 \(\gamma\) 固定主尺度，结构圆维 \(\cdim\) 固定可见连续相位通道的结构维数。
合并即得等价性与完备性。证毕。
\end{proof}

\begin{conjecture}[混合型纤维谱的绝对连续性与二阶曲线光滑性]\label{conj:xi-zero-dispersion-atomless-smooth-b}
在满足强混合与非退化 Gibbs 结构的读出因子映射类中，缩放极限 \(Z\) 无原子且具有 \(C^{1}\) 密度；
等价地，\(\varepsilon\mapsto b(\varepsilon)\) 在 \((0,1)\) 上无折点并且为 \(C^{1}\) 函数。
有限原子零色散对应高度刚性的退化谱型；一旦存在真正的多尺度纤维涨落，原子结构将被连续谱所取代。
\end{conjecture}

\par
综上，有限原子零色散极限将常数偏置成功率曲线压缩为分段仿射对象，并使二阶常数项 \(b(\varepsilon)\) 具有闭式反演与折点计数；全曲线的可识别性把 \(b(\varepsilon)\) 提升为可复核的分布律接口；而串接不等式则给出审计链条组合时二阶预算的结构上界。
与圆维签名 \(\Sigma\Sigma\) 一并使用时，连续圆因子与离散残差轴在同一制度内获得清晰分工：结构层由 \(\Sigma_{\circ}\) 计量，二阶实现代价由 \(\gamma,b(\cdot)\) 量化并由 \(\Sigma_{\mathrm{res}}\) 外置封口。

\endinput
